\documentclass[11pt,a4paper]{report}
% Compilation : tectonic Cours_Optimisation_Bayesienne.tex
% Figures à données : python tools/make_figures.py (lit Code/resultats/pipeline_L1_V3.json et diagnostic_L1_V3.json,
% produits par python -m optim.experiences_l1 et python -m optim.diagnostic_l1)
% Préambule commun (repris du cours d'optimisation bi-objectif)
\usepackage{iftex}
\ifPDFTeX
  \usepackage[utf8]{inputenc}
  \usepackage[T1]{fontenc}
  \usepackage{lmodern}
\else
  \usepackage{fontspec}
\fi
\usepackage[french]{babel}
\usepackage[a4paper,margin=2.3cm]{geometry}
\usepackage{amsmath,amssymb,mathtools}
\usepackage{xcolor}
\usepackage{tikz}
\usetikzlibrary{arrows.meta,calc,positioning,shapes.geometric,decorations.pathreplacing,decorations.markings,patterns,babel}
\usepackage{pgfplots}
\pgfplotsset{compat=1.18}
\usepgfplotslibrary{groupplots}
\usepackage{tabularx,booktabs,array}
\usepackage{enumitem}
\setlist{nosep,leftmargin=*}
\usepackage{listings}
\usepackage[font=small,labelfont=bf]{caption}
\usepackage[raggedright]{titlesec}
\usepackage[most]{tcolorbox}
\usepackage[colorlinks=true,linkcolor=blue!50!black,urlcolor=blue!50!black,citecolor=blue!50!black]{hyperref}

\definecolor{cblue}{HTML}{2A78D6}
\definecolor{corange}{HTML}{EB6834}
\definecolor{cgreen}{HTML}{1BAF7A}
\definecolor{cink}{HTML}{1F2328}
\definecolor{cmuted}{HTML}{6B6F76}

% ---------------------------------------------------------------- notations
\newcommand{\dd}{\mathrm{d}}
\newcommand{\Ue}{U_e}
\renewcommand{\Re}{\mathit{Re}}
\newcommand{\vect}[1]{\mathbf{#1}}
\newcommand{\code}[1]{\texttt{#1}}

% ---------------------------------------------------------------- encadrés
\tcbset{coursbox/.style={enhanced,breakable,boxrule=0pt,leftrule=2.5pt,arc=0pt,
  left=6pt,right=6pt,top=4pt,bottom=4pt,fonttitle=\bfseries\small,coltitle=cink,
  attach title to upper={\ }}}
\newtcolorbox{definition}[1][]{coursbox,colback=cblue!5,colframe=cblue,title={Définition#1.}}
\newtcolorbox{resultat}[1][]{coursbox,colback=cgreen!6,colframe=cgreen,title={À retenir#1.}}
\newtcolorbox{projet}[1][]{coursbox,colback=corange!6,colframe=corange,title={Dans notre projet#1.}}
\newtcolorbox{piege}[1][]{coursbox,colback=cmuted!7,colframe=cmuted,title={Piège#1.}}
\newtcolorbox{codebox}[1][]{coursbox,colback=cmuted!3,colframe=cblue!60,title={Dans le code#1.}}

% ---------------------------------------------------------------- exercices
\newcounter{exo}[chapter]
\renewcommand{\theexo}{\thechapter.\arabic{exo}}
\newenvironment{exercice}[1][]{\refstepcounter{exo}\par\medskip\noindent
  \textbf{\color{cblue}Exercice~\theexo}\if\relax\detokenize{#1}\relax\else\ (#1)\fi.\ }{\par}
\newcommand{\exercices}{\section*{Exercices}\addcontentsline{toc}{section}{Exercices}
  \noindent\textit{Corrigés en annexe~\ref{annexe:corriges}. Chercher d'abord sans les regarder.}}

\tikzset{
  flow/.style={draw=cink,rounded corners=2pt,fill=white,align=center,font=\small,inner sep=4pt},
  arr/.style={-{Stealth[length=2.2mm]},cink,thick},
}

\lstset{basicstyle=\ttfamily\footnotesize,breaklines=true,columns=fullflexible,frame=single,rulecolor=\color{cmuted!50},
  backgroundcolor=\color{cmuted!6},keywordstyle=\color{cblue}\bfseries,commentstyle=\color{cmuted}\itshape,
  stringstyle=\color{corange},showstringspaces=false,language=Python,
  literate={é}{\'e}1{è}{\`e}1{à}{\`a}1{ê}{\^e}1{ç}{\c c}1{ô}{\^o}1{î}{\^i}1{ù}{\`u}1{É}{\'E}1{→}{$\to$}1}

\setcounter{secnumdepth}{2}
\setcounter{tocdepth}{1}
\setlength{\parindent}{0pt}
\setlength{\emergencystretch}{3em}
\setlength{\parskip}{0.45em}


\definecolor{cviolet}{HTML}{7A4FB3}
\newtcolorbox{pistes}[1][]{coursbox,colback=cviolet!6,colframe=cviolet,title={Pistes d'amélioration#1.}}
\newcommand{\dB}{\,\mathrm{dB}}


% ============================================================================
% Itérations du projet HackathonQuantique : couleurs et marqueurs communs.
% Inclus par tous les cours et par la fiche projet, pour qu'une itération ait
% la même couleur partout. Pour ajouter une itération : une couleur, un nom,
% une date, une ligne dans \legendeiterations.
% ============================================================================
\definecolor{itV0}{HTML}{8D6E63}   % brun   : V0, modèle de la note d'intention
\definecolor{itV1}{HTML}{C2185B}   % framboise : V1, chaîne complète conforme à la note
\definecolor{itV2}{HTML}{00838F}   % cyan foncé : V2, problème reformulé
\definecolor{itV3}{HTML}{E09A00}   % ambre  : V3, concurrent équitable et problème plus grand
\definecolor{itV4}{HTML}{5E35B1}   % indigo : V4, réservé

\makeatletter
\newcommand{\itername}[1]{\@nameuse{itname@#1}}
\newcommand{\iterdate}[1]{\@nameuse{itdate@#1}}
\newcommand{\defiter}[3]{\@namedef{itname@#1}{#2}\@namedef{itdate@#1}{#3}}
\makeatother
\defiter{V0}{Modèle de la note d'intention}{24--30 oct. 2026}
\defiter{V1}{Chaîne complète conforme à la note}{4--5 nov. 2026}
\defiter{V2}{Problème reformulé}{5--6 nov. 2026}
\defiter{V3}{Concurrent équitable, problème plus grand}{à partir du 6 nov. 2026}

% Pastille : \iterbadge{V1}
\newcommand{\iterbadge}[1]{\tikz[baseline=(b.base)]\node[fill=it#1,text=white,rounded corners=2pt,
  inner xsep=3pt,inner ysep=1.5pt,font=\sffamily\bfseries\scriptsize] (b) {#1};}

% Bandeau de début de chapitre : \iterbanner{V0}
\newcommand{\iterbanner}[1]{\par\noindent\begin{tikzpicture}
  \node[fill=it#1!12,draw=it#1,line width=0.6pt,rounded corners=2pt,text width=\dimexpr\linewidth-12pt\relax,
    inner sep=4pt,anchor=west,font=\small] at (0,0)
    {\iterbadge{#1}\ \textbf{\color{it#1}Itération #1} --- \itername{#1} \hfill{\color{it#1}\iterdate{#1}}};
  \end{tikzpicture}\par\medskip}

% Encadré d'évolution : \begin{evolution}{V1}{titre} ... \end{evolution}
\newtcolorbox{evolution}[2]{enhanced,breakable,boxrule=0pt,leftrule=3pt,arc=0pt,
  left=6pt,right=6pt,top=4pt,bottom=4pt,colback=it#1!7,colframe=it#1,
  fonttitle=\bfseries\small,coltitle=it#1,attach title to upper={\ },
  title={\iterbadge{#1}\ Évolution #1 : #2.}}

% Légende des itérations, pour les avant-propos
\newcommand{\legendeiterations}{%
\begin{center}\small
\begin{tabular}{lll}
\iterbadge{V0} & \itername{V0} & \iterdate{V0}\\
\iterbadge{V1} & \itername{V1} & \iterdate{V1}\\
\iterbadge{V2} & \itername{V2} & \iterdate{V2}\\
\iterbadge{V3} & \itername{V3} & \iterdate{V3}\\
\end{tabular}
\end{center}}

% Encadré « pistes d'amélioration » : défini ici s'il manque (cours V0 sans cet encadré)
\providecolor{cviolet}{HTML}{7A4FB3}
\ifcsname pistes\endcsname\else
\newtcolorbox{pistes}[1][]{enhanced,breakable,boxrule=0pt,leftrule=2.5pt,arc=0pt,
  left=6pt,right=6pt,top=4pt,bottom=4pt,fonttitle=\bfseries\small,coltitle=cink,attach title to upper={\ },
  colback=cviolet!6,colframe=cviolet,title={Pistes d'amélioration#1.}}
\fi


\title{\textbf{Optimisation bayésienne}\\[0.4em]\large Krigeage, amélioration espérée et contraintes :\\ le concurrent de même nature de la boucle QUBO (lot L1 de la V3)}
\author{Notes de travail, projet \emph{HackathonQuantique}}
\date{Version du 15 novembre 2026}

\begin{document}
\pagenumbering{Alph}
\maketitle
\pagenumbering{roman}

\chapter*{Avant-propos}
\addcontentsline{toc}{chapter}{Avant-propos}
\paragraph{Objet.} La V2 a montré que la boucle QUBO retrouve le front de Pareto traînée/SER quand NSGA-II n'y parvient presque jamais. Mais NSGA-II n'utilise aucun modèle, alors que la boucle en ajuste un à chaque itération. Le concurrent de même nature est l'\textbf{optimisation bayésienne}. Ce cours la construit depuis le début, l'applique au problème V2 avec le même protocole, et répond à l'hypothèse H1 de la fiche projet : \emph{l'optimisation bayésienne égale-t-elle ou bat-elle la boucle QUBO ?}

\paragraph{Ce qu'on suppose connu.} Le cours d'optimisation bi-objectif (front de Pareto, hypervolume, boucle QUBO, protocole de comparaison) et son chapitre V2. Côté mathématiques : la loi normale, les moindres carrés, un peu d'algèbre linéaire. Le krigeage et l'amélioration espérée sont posés depuis zéro.

\paragraph{Fil conducteur.} La question et le plan d'expérience (ch.~\ref{ch:question}) $\to$ le krigeage, un modèle qui donne son incertitude (ch.~\ref{ch:krigeage}) $\to$ l'amélioration espérée, qui s'en sert pour choisir le prochain profil (ch.~\ref{ch:acquisition}) $\to$ contraintes et deux objectifs (ch.~\ref{ch:contraintes}) $\to$ le code et sa validation (ch.~\ref{ch:code}) $\to$ les résultats de la comparaison (ch.~\ref{ch:resultats}) $\to$ ce qui les explique : la contrainte de manœuvre (ch.~\ref{ch:diagnostic}) $\to$ bilan, limites et pistes (ch.~\ref{ch:bilan}).

\paragraph{Figures et chiffres.} Les courbes sont calculées par le code du projet (\code{Code/optim/bayes.py}) via \code{tools/make\_figures.py}, à partir de \code{resultats/pipeline\_L1\_V3.json} (produit par \code{python -m optim.experiences\_l1}) et de \code{resultats/diagnostic\_L1\_V3.json} (produit par \code{python -m optim.diagnostic\_l1}). Chaque chapitre se termine par des exercices (corrigés en annexe~\ref{annexe:corriges}) et des \textbf{pistes d'amélioration}.

\paragraph{Itérations du projet.} Tout ce cours appartient à l'itération V3 (lot L1) : chaque chapitre porte son bandeau.
\legendeiterations

\tableofcontents
\cleardoublepage
\pagenumbering{arabic}

\chapter{La question : quel est le vrai concurrent de la boucle QUBO ?}\label{ch:question}
\iterbanner{V3}

\section{Ce que la V2 a laissé ouvert}
Sur le problème V2 (1\,024 profils, deux objectifs, contrainte de manœuvre), la boucle QUBO à solveur exact retrouve tout le front exact en 75 simulations dans le scénario S1 et 142 dans S2 (médianes sur 20 graines). NSGA-II n'y parvient presque jamais dans un budget de 250. La fiche projet en tire une réserve : cette comparaison oppose deux méthodes qui ne sont pas de même nature.

\begin{definition}[ : méthode à modèle de substitution]
Une méthode d'optimisation est dite \emph{à modèle de substitution} si, à chaque itération, elle (1) ajuste un modèle peu coûteux de la fonction à optimiser sur les points déjà simulés, puis (2) se sert de ce modèle pour choisir les prochains points à simuler.
\end{definition}

NSGA-II ne fait ni l'un ni l'autre : il recombine les meilleurs profils connus. La boucle QUBO, elle, est une méthode à modèle de substitution. Quand elle bat NSGA-II, on ne sait donc pas ce qui a gagné : la formulation QUBO, ou le simple fait d'avoir un modèle.

\begin{projet}
L'optimisation bayésienne est la méthode à modèle de substitution de référence pour les fonctions coûteuses, en particulier en conception aéronautique (krigeage, algorithme EGO, outils de l'ONERA et de l'ISAE comme la bibliothèque SMT). C'est le concurrent qu'un jury attend.
\end{projet}

\section{Deux choix indépendants}
Toute méthode à modèle de substitution fait deux choix, qu'on peut séparer.

\begin{center}
\renewcommand{\arraystretch}{1.25}
\begin{tabularx}{\linewidth}{lXX}
\toprule
 & Boucle QUBO (V2) & Optimisation bayésienne \\
\midrule
\textbf{Modèle} & polynôme quadratique ajusté par moindres carrés & krigeage (processus gaussien) : une prédiction \emph{et} une incertitude \\
\textbf{Sélection} & minimum du modèle, trouvé en résolvant un QUBO & maximum d'un \emph{critère d'acquisition} qui pèse prédiction et incertitude \\
\bottomrule
\end{tabularx}
\end{center}

\begin{figure}[htbp]
\centering
\begin{tikzpicture}[node distance=0.5cm and 0.6cm]
\node[flow] (init) {10 profils\\ tirés au hasard};
\node[flow,right=of init,fill=cblue!8] (fit) {\textbf{modèle}\\ quadratique\\ ou krigeage};
\node[flow,right=of fit,fill=corange!8] (sel) {\textbf{sélection}, pour chaque poids $w$ :\\ minimum de la prédiction\\ ou amélioration espérée};
\node[flow,below=1.0cm of sel] (sim) {simuler les 5 profils\\ choisis (un par poids)};
\node[flow,below=1.0cm of fit] (arch) {archive des\\ profils simulés};
\draw[arr] (init) -- (fit); \draw[arr] (fit) -- (sel);
\draw[arr] (sel) -- (sim); \draw[arr] (sim) -- (arch); \draw[arr] (arch) -- (fit);
\end{tikzpicture}
\caption{Le schéma commun. La boucle QUBO et l'optimisation bayésienne ne diffèrent que par les deux cases colorées.}\label{fig:schema}
\end{figure}

\section{Le plan d'expérience : une ablation}
Pour savoir lequel des deux choix compte, on les croise. C'est une \emph{ablation} : on ne change qu'un élément à la fois.

\begin{center}
\renewcommand{\arraystretch}{1.25}
\begin{tabularx}{\linewidth}{l>{\centering\arraybackslash}X>{\centering\arraybackslash}X}
\toprule
 & sélection : minimum de la prédiction & sélection : amélioration espérée \\
\midrule
modèle quadratique & contrôle (la boucle QUBO, sans encodage binaire) & effet de la sélection seule \\
krigeage & effet du modèle seul & optimisation bayésienne \\
\bottomrule
\end{tabularx}
\end{center}

La case en haut à gauche sert de \textbf{contrôle} : elle fait la même chose que la boucle QUBO résolue exactement, mais par simple énumération de la grille, sans bits ni qubits auxiliaires. Si elle se comporte comme la boucle QUBO, les trois autres cases sont interprétables.

\begin{resultat}[ : l'hypothèse testée]
\textbf{H1} (fiche projet, V3) : \emph{sur le problème V2, l'optimisation bayésienne égale ou bat la boucle QUBO, parce que son modèle (un processus gaussien) est plus souple qu'un modèle quadratique.}
\end{resultat}

\section{Ce qui reste identique}
Pour que la comparaison soit équitable, tout le reste est repris de la V2 (\code{optim/experiences\_v2.py}) :
\begin{itemize}
\item la table XFOIL des 1\,024 profils et la SER par méthode des moments ;
\item le compteur de simulations \code{CountedProblem} : un profil déjà simulé n'est pas recompté ;
\item les 10 profils initiaux, \emph{les mêmes} que ceux de la boucle QUBO pour une graine donnée ;
\item les 5 poids $w\in\{0;\ 0{,}25;\ 0{,}5;\ 0{,}75;\ 1\}$ et un profil par poids et par itération ;
\item le budget de 250 simulations, 20 graines, et les deux critères d'arrêt (front exact retrouvé, arrêt réaliste par stagnation de l'hypervolume).
\end{itemize}
Comme les graines et les points initiaux sont communs, les exécutions sont \emph{appariées} : on peut comparer deux méthodes graine par graine (chapitre~\ref{ch:resultats}).

\exercices
\begin{exercice}
Classer ces méthodes en « à modèle de substitution » ou non : tirage aléatoire, NSGA-II, recuit simulé appliqué directement aux simulations, boucle QUBO, optimisation bayésienne. Pour celles qui en ont un, dire quel est le modèle.
\end{exercice}
\begin{exercice}
Dans l'ablation, pourquoi ne suffit-il pas de comparer directement la boucle QUBO à l'optimisation bayésienne (les deux cases de la diagonale) ?
\end{exercice}

\begin{pistes}
\begin{itemize}
\item Ajouter à l'ablation un troisième modèle (forêt aléatoire, réseau de neurones) pour voir si la conclusion sur le modèle se généralise.
\item Refaire l'ablation sur la grande grille du lot L2 (16\,384 profils), où l'énumération de l'acquisition n'est plus gratuite.
\end{itemize}
\end{pistes}

\chapter{Le krigeage : un modèle qui donne son incertitude}\label{ch:krigeage}
\iterbanner{V3}

\section{Ce qui manque au modèle quadratique}
Le modèle quadratique de la boucle QUBO rend une valeur par profil, rien d'autre. Il ne dit pas s'il prédit un profil voisin de dix profils déjà simulés, ou un profil isolé dans un coin de la grille. Or cette information compte : là où le modèle est sûr de lui, simuler n'apprend rien ; là où il ne sait pas, un bon profil peut se cacher.

Le \textbf{krigeage}, ou \emph{régression par processus gaussien}, rend pour chaque profil une prédiction $\mu(x)$ et un écart-type $s(x)$ qui mesure l'incertitude de cette prédiction.

\section{Processus gaussien}
\begin{definition}[ : processus gaussien]
Un processus gaussien est une loi de probabilité sur des fonctions $f$, telle que pour tout ensemble fini de points $x_1,\dots,x_n$ le vecteur $(f(x_1),\dots,f(x_n))$ suit une loi normale multivariée. Il est entièrement décrit par une moyenne $m(x)$ et un \emph{noyau} (ou fonction de covariance)
\[
k(x,x')=\mathrm{Cov}\big(f(x),f(x')\big).
\]
\end{definition}

Le noyau encode l'hypothèse de régularité : deux profils proches ont des valeurs corrélées, deux profils éloignés n'en ont presque pas. On utilise ici le noyau \textbf{Matérn 5/2}, produit sur les $d$ variables de conception :
\[
k(x,x')=\sigma^2\prod_{i=1}^{d}\Big(1+\sqrt5\,\theta_i\,|x_i-x_i'|+\tfrac53\,\theta_i^2\,(x_i-x_i')^2\Big)\exp\!\big(-\sqrt5\,\theta_i\,|x_i-x_i'|\big).
\]
$\sigma^2$ est la variance du processus. Chaque $\theta_i>0$ règle la vitesse à laquelle la corrélation décroît le long de la variable $i$ : un grand $\theta_i$ signifie que la fonction varie vite selon cette variable. Les fonctions tirées d'un Matérn 5/2 sont deux fois dérivables, ce qui est moins exigeant que le noyau gaussien (infiniment dérivable) et convient mieux à des sorties de simulation.

\section{Conditionner sur les simulations}
On a simulé $n$ profils : $X=(x_1,\dots,x_n)$, valeurs $\vect y=(y_1,\dots,y_n)$. On suppose $y_j=f(x_j)+\varepsilon_j$, avec un bruit $\varepsilon_j\sim\mathcal N(0,\sigma_n^2)$ indépendant. Notons
\[
K=\big[k(x_j,x_l)\big]_{j,l}+\sigma_n^2 I,\qquad \vect k_\star(x)=\big[k(x_j,x)\big]_j .
\]
Avec une moyenne constante $m$, la loi de $f(x)$ sachant les simulations est normale, de moyenne et de variance
\begin{equation}\label{eq:krigeage}
\mu(x)=m+\vect k_\star(x)^{\!\top}K^{-1}(\vect y-m\vect 1),\qquad
s^2(x)=k(x,x)-\vect k_\star(x)^{\!\top}K^{-1}\vect k_\star(x).
\end{equation}

\begin{resultat}
\begin{itemize}
\item $\mu(x)$ est une combinaison linéaire des valeurs simulées, pondérée par la corrélation avec $x$.
\item $s^2(x)$ part de la variance a priori $k(x,x)=\sigma^2$ et \emph{diminue} de ce que les simulations voisines apprennent sur $x$. Elle ne dépend pas des valeurs $\vect y$, seulement de la position des points.
\item Sans bruit ($\sigma_n=0$), le krigeage interpole : $\mu(x_j)=y_j$ et $s(x_j)=0$.
\end{itemize}
\end{resultat}

\begin{figure}[htbp]
\centering
% figure générée par tools/make_figures.py : gp_1d
\begin{tikzpicture}
\begin{groupplot}[group style={group size=1 by 2,vertical sep=0.5cm,xticklabels at=edge bottom},width=0.82\linewidth,
  xmin=-1,xmax=1,tick label style={font=\scriptsize},label style={font=\small},
  legend style={draw=none,font=\scriptsize,at={(0.5,1.03)},anchor=south},legend columns=4]
\nextgroupplot[height=5.6cm,ylabel={$f(x)$}]
\addplot[fill=cblue!15,draw=none,forget plot] coordinates {(-1,1.1925) (-0.9875,1.2279) (-0.975,1.2772) (-0.9625,1.3242) (-0.95,1.3675) (-0.9375,1.4068) (-0.925,1.4419) (-0.9125,1.4729) (-0.9,1.5) (-0.8875,1.5235) (-0.875,1.5434) (-0.8625,1.5601) (-0.85,1.5738) (-0.8375,1.5849) (-0.825,1.5936) (-0.8125,1.6001) (-0.8,1.6047) (-0.7875,1.6076) (-0.775,1.609) (-0.7625,1.6092) (-0.75,1.6083) (-0.7375,1.6065) (-0.725,1.604) (-0.7125,1.6007) (-0.7,1.597) (-0.6875,1.5928) (-0.675,1.5883) (-0.6625,1.5835) (-0.65,1.5786) (-0.6375,1.5734) (-0.625,1.5682) (-0.6125,1.5629) (-0.6,1.5575) (-0.5875,1.5521) (-0.575,1.5467) (-0.5625,1.5412) (-0.55,1.5357) (-0.5375,1.5302) (-0.525,1.5246) (-0.5125,1.519) (-0.5,1.5133) (-0.4875,1.5075) (-0.475,1.5015) (-0.4625,1.4955) (-0.45,1.4892) (-0.4375,1.4827) (-0.425,1.476) (-0.4125,1.469) (-0.4,1.4616) (-0.3875,1.4539) (-0.375,1.4457) (-0.3625,1.4371) (-0.35,1.4279) (-0.3375,1.4182) (-0.325,1.4078) (-0.3125,1.3967) (-0.3,1.3848) (-0.2875,1.3721) (-0.275,1.3584) (-0.2625,1.3438) (-0.25,1.3281) (-0.2375,1.3113) (-0.225,1.2932) (-0.2125,1.2738) (-0.2,1.2531) (-0.1875,1.2308) (-0.175,1.207) (-0.1625,1.1816) (-0.15,1.1544) (-0.1375,1.1255) (-0.125,1.0948) (-0.1125,1.0622) (-0.1,1.0277) (-0.0875,0.99138) (-0.075,0.95316) (-0.0625,0.91313) (-0.05,0.8714) (-0.0375,0.82808) (-0.025,0.78339) (-0.0125,0.73756) (0,0.6909) (0.0125,0.64381) (0.025,0.59677) (0.0375,0.55035) (0.05,0.50524) (0.0625,0.46228) (0.075,0.42265) (0.0875,0.38869) (0.1,0.37479) (0.1125,0.41622) (0.125,0.47195) (0.1375,0.52686) (0.15,0.57839) (0.1625,0.62522) (0.175,0.66636) (0.1875,0.70104) (0.2,0.72863) (0.2125,0.74867) (0.225,0.76089) (0.2375,0.76519) (0.25,0.76169) (0.2625,0.75081) (0.275,0.73346) (0.2875,0.71202) (0.3,0.69916) (0.3125,0.72053) (0.325,0.74985) (0.3375,0.7742) (0.35,0.791) (0.3625,0.79917) (0.375,0.79818) (0.3875,0.78775) (0.4,0.7679) (0.4125,0.73885) (0.425,0.70103) (0.4375,0.65508) (0.45,0.60188) (0.4625,0.54259) (0.475,0.47883) (0.4875,0.41366) (0.5,0.36018) (0.5125,0.34379) (0.525,0.33966) (0.5375,0.33527) (0.55,0.32815) (0.5625,0.31725) (0.575,0.30206) (0.5875,0.28233) (0.6,0.25807) (0.6125,0.22953) (0.625,0.19715) (0.6375,0.16161) (0.65,0.12381) (0.6625,0.084911) (0.675,0.046513) (0.6875,0.011597) (0.7,-0.0066895) (0.7125,0.017902) (0.725,0.059401) (0.7375,0.10456) (0.75,0.15022) (0.7625,0.1946) (0.775,0.23631) (0.7875,0.2742) (0.8,0.30733) (0.8125,0.33494) (0.825,0.35645) (0.8375,0.3715) (0.85,0.37991) (0.8625,0.38179) (0.875,0.37766) (0.8875,0.3694) (0.9,0.37157) (0.9125,0.42481) (0.925,0.49368) (0.9375,0.56375) (0.95,0.63278) (0.9625,0.69978) (0.975,0.76417) (0.9875,0.82554) (1,0.88363) (1,0.11645) (0.9875,0.14739) (0.975,0.17862) (0.9625,0.20971) (0.95,0.24011) (0.9375,0.26912) (0.925,0.29571) (0.9125,0.31765) (0.9,0.32064) (0.8875,0.26946) (0.875,0.20515) (0.8625,0.14276) (0.85,0.084805) (0.8375,0.032508) (0.825,-0.013233) (0.8125,-0.051723) (0.8,-0.08242) (0.7875,-0.10493) (0.775,-0.11903) (0.7625,-0.12466) (0.75,-0.12196) (0.7375,-0.11137) (0.725,-0.093792) (0.7125,-0.071636) (0.7,-0.057587) (0.6875,-0.077327) (0.675,-0.10464) (0.6625,-0.1267) (0.65,-0.141) (0.6375,-0.14658) (0.625,-0.14298) (0.6125,-0.13003) (0.6,-0.1078) (0.5875,-0.076578) (0.575,-0.036837) (0.5625,0.010744) (0.55,0.065273) (0.5375,0.12559) (0.525,0.19007) (0.5125,0.25567) (0.5,0.30929) (0.4875,0.32554) (0.475,0.32924) (0.4625,0.33291) (0.45,0.33901) (0.4375,0.34857) (0.425,0.36213) (0.4125,0.37994) (0.4,0.40203) (0.3875,0.4282) (0.375,0.45805) (0.3625,0.49099) (0.35,0.52619) (0.3375,0.56259) (0.325,0.59869) (0.3125,0.6316) (0.3,0.64826) (0.2875,0.62248) (0.275,0.58027) (0.2625,0.53489) (0.25,0.48951) (0.2375,0.44594) (0.225,0.40556) (0.2125,0.36955) (0.2,0.33889) (0.1875,0.31439) (0.175,0.2967) (0.1625,0.28625) (0.15,0.28331) (0.1375,0.28785) (0.125,0.29945) (0.1125,0.31629) (0.1,0.32386) (0.0875,0.28154) (0.075,0.22471) (0.0625,0.16771) (0.05,0.11266) (0.0375,0.0604) (0.025,0.011391) (0.0125,-0.034112) (0,-0.075985) (-0.0125,-0.1142) (-0.025,-0.14882) (-0.0375,-0.17994) (-0.05,-0.2077) (-0.0625,-0.2323) (-0.075,-0.25392) (-0.0875,-0.27278) (-0.1,-0.28908) (-0.1125,-0.30304) (-0.125,-0.31487) (-0.1375,-0.32477) (-0.15,-0.33292) (-0.1625,-0.33952) (-0.175,-0.34472) (-0.1875,-0.34867) (-0.2,-0.35153) (-0.2125,-0.35342) (-0.225,-0.35446) (-0.2375,-0.35473) (-0.25,-0.35434) (-0.2625,-0.35337) (-0.275,-0.35187) (-0.2875,-0.34991) (-0.3,-0.34754) (-0.3125,-0.34479) (-0.325,-0.34169) (-0.3375,-0.33826) (-0.35,-0.33453) (-0.3625,-0.3305) (-0.375,-0.32616) (-0.3875,-0.32153) (-0.4,-0.31659) (-0.4125,-0.31133) (-0.425,-0.30574) (-0.4375,-0.29978) (-0.45,-0.29343) (-0.4625,-0.28667) (-0.475,-0.27946) (-0.4875,-0.27176) (-0.5,-0.26353) (-0.5125,-0.25473) (-0.525,-0.2453) (-0.5375,-0.2352) (-0.55,-0.22436) (-0.5625,-0.21274) (-0.575,-0.20025) (-0.5875,-0.18685) (-0.6,-0.17245) (-0.6125,-0.15698) (-0.625,-0.14036) (-0.6375,-0.12252) (-0.65,-0.10337) (-0.6625,-0.08282) (-0.675,-0.060781) (-0.6875,-0.037165) (-0.7,-0.011878) (-0.7125,0.015171) (-0.725,0.044073) (-0.7375,0.074918) (-0.75,0.10779) (-0.7625,0.14277) (-0.775,0.17993) (-0.7875,0.21933) (-0.8,0.26102) (-0.8125,0.30503) (-0.825,0.35138) (-0.8375,0.40005) (-0.85,0.45101) (-0.8625,0.50417) (-0.875,0.55943) (-0.8875,0.61663) (-0.9,0.67556) (-0.9125,0.73596) (-0.925,0.7975) (-0.9375,0.85978) (-0.95,0.92231) (-0.9625,0.98448) (-0.975,1.0454) (-0.9875,1.1033) (-1,1.1416)};
\addplot[cmuted,thick,dashed] coordinates {(-1,1.1676) (-0.9875,1.1855) (-0.975,1.2025) (-0.9625,1.2184) (-0.95,1.2329) (-0.9375,1.2459) (-0.925,1.2572) (-0.9125,1.2665) (-0.9,1.2737) (-0.8875,1.2785) (-0.875,1.281) (-0.8625,1.2808) (-0.85,1.278) (-0.8375,1.2723) (-0.825,1.2638) (-0.8125,1.2522) (-0.8,1.2377) (-0.7875,1.2201) (-0.775,1.1995) (-0.7625,1.1758) (-0.75,1.149) (-0.7375,1.1193) (-0.725,1.0867) (-0.7125,1.0512) (-0.7,1.0129) (-0.6875,0.97209) (-0.675,0.92874) (-0.6625,0.88304) (-0.65,0.83516) (-0.6375,0.78527) (-0.625,0.73356) (-0.6125,0.68024) (-0.6,0.62551) (-0.5875,0.56961) (-0.575,0.51275) (-0.5625,0.45518) (-0.55,0.39715) (-0.5375,0.33889) (-0.525,0.28067) (-0.5125,0.22273) (-0.5,0.16533) (-0.4875,0.10871) (-0.475,0.053138) (-0.4625,-0.0011557) (-0.45,-0.053928) (-0.4375,-0.10495) (-0.425,-0.15399) (-0.4125,-0.20083) (-0.4,-0.24528) (-0.3875,-0.28713) (-0.375,-0.32622) (-0.3625,-0.36237) (-0.35,-0.39543) (-0.3375,-0.42526) (-0.325,-0.45175) (-0.3125,-0.4748) (-0.3,-0.49431) (-0.2875,-0.51022) (-0.275,-0.52249) (-0.2625,-0.53109) (-0.25,-0.536) (-0.2375,-0.53723) (-0.225,-0.53481) (-0.2125,-0.52879) (-0.2,-0.51922) (-0.1875,-0.5062) (-0.175,-0.48983) (-0.1625,-0.47021) (-0.15,-0.4475) (-0.1375,-0.42182) (-0.125,-0.39336) (-0.1125,-0.36228) (-0.1,-0.32879) (-0.0875,-0.29307) (-0.075,-0.25535) (-0.0625,-0.21586) (-0.05,-0.17481) (-0.0375,-0.13246) (-0.025,-0.089038) (-0.0125,-0.044802) (0,0) (0.0125,0.045114) (0.025,0.090288) (0.0375,0.13527) (0.05,0.17981) (0.0625,0.22367) (0.075,0.2666) (0.0875,0.30838) (0.1,0.34879) (0.1125,0.38759) (0.125,0.42461) (0.1375,0.45963) (0.15,0.4925) (0.1625,0.52303) (0.175,0.55108) (0.1875,0.57652) (0.2,0.59922) (0.2125,0.6191) (0.225,0.63606) (0.2375,0.65004) (0.25,0.661) (0.2625,0.6689) (0.275,0.67374) (0.2875,0.67554) (0.3,0.67431) (0.3125,0.67011) (0.325,0.663) (0.3375,0.65307) (0.35,0.64043) (0.3625,0.62518) (0.375,0.60747) (0.3875,0.58745) (0.4,0.56528) (0.4125,0.54114) (0.425,0.51524) (0.4375,0.48776) (0.45,0.45893) (0.4625,0.42897) (0.475,0.39811) (0.4875,0.3666) (0.5,0.33467) (0.5125,0.30258) (0.525,0.27058) (0.5375,0.23892) (0.55,0.20785) (0.5625,0.17763) (0.575,0.1485) (0.5875,0.12071) (0.6,0.094488) (0.6125,0.070074) (0.625,0.047688) (0.6375,0.027543) (0.65,0.0098403) (0.6625,-0.0052297) (0.675,-0.01749) (0.6875,-0.026778) (0.7,-0.032945) (0.7125,-0.03586) (0.725,-0.035407) (0.7375,-0.031486) (0.75,-0.024018) (0.7625,-0.01294) (0.775,0.0017923) (0.7875,0.020204) (0.8,0.042301) (0.8125,0.068071) (0.825,0.097483) (0.8375,0.13049) (0.85,0.16701) (0.8625,0.20697) (0.875,0.25026) (0.8875,0.29677) (0.9,0.34634) (0.9125,0.39884) (0.925,0.45408) (0.9375,0.5119) (0.95,0.57209) (0.9625,0.63445) (0.975,0.69877) (0.9875,0.76481) (1,0.83235)};\addlegendentry{fonction}
\addplot[cblue,very thick] coordinates {(-1,1.1671) (-0.9875,1.1656) (-0.975,1.1613) (-0.9625,1.1543) (-0.95,1.1449) (-0.9375,1.1333) (-0.925,1.1197) (-0.9125,1.1044) (-0.9,1.0878) (-0.8875,1.07) (-0.875,1.0514) (-0.8625,1.0321) (-0.85,1.0124) (-0.8375,0.99248) (-0.825,0.97248) (-0.8125,0.95255) (-0.8,0.93285) (-0.7875,0.91346) (-0.775,0.89449) (-0.7625,0.876) (-0.75,0.85807) (-0.7375,0.84073) (-0.725,0.82402) (-0.7125,0.80795) (-0.7,0.79256) (-0.6875,0.77783) (-0.675,0.76377) (-0.6625,0.75036) (-0.65,0.7376) (-0.6375,0.72546) (-0.625,0.71392) (-0.6125,0.70295) (-0.6,0.69253) (-0.5875,0.68262) (-0.575,0.6732) (-0.5625,0.66423) (-0.55,0.65566) (-0.5375,0.64748) (-0.525,0.63965) (-0.5125,0.63212) (-0.5,0.62487) (-0.4875,0.61785) (-0.475,0.61104) (-0.4625,0.60439) (-0.45,0.59788) (-0.4375,0.59147) (-0.425,0.58513) (-0.4125,0.57882) (-0.4,0.57251) (-0.3875,0.56617) (-0.375,0.55978) (-0.3625,0.55329) (-0.35,0.54669) (-0.3375,0.53995) (-0.325,0.53304) (-0.3125,0.52594) (-0.3,0.51862) (-0.2875,0.51108) (-0.275,0.50328) (-0.2625,0.49522) (-0.25,0.48689) (-0.2375,0.47827) (-0.225,0.46938) (-0.2125,0.46021) (-0.2,0.45077) (-0.1875,0.44107) (-0.175,0.43115) (-0.1625,0.42104) (-0.15,0.41076) (-0.1375,0.40039) (-0.125,0.38997) (-0.1125,0.37959) (-0.1,0.36933) (-0.0875,0.3593) (-0.075,0.34962) (-0.0625,0.34042) (-0.05,0.33185) (-0.0375,0.32407) (-0.025,0.31729) (-0.0125,0.31168) (0,0.30746) (0.0125,0.30485) (0.025,0.30408) (0.0375,0.30537) (0.05,0.30895) (0.0625,0.31499) (0.075,0.32368) (0.0875,0.33511) (0.1,0.34933) (0.1125,0.36625) (0.125,0.3857) (0.1375,0.40736) (0.15,0.43085) (0.1625,0.45574) (0.175,0.48153) (0.1875,0.50772) (0.2,0.53376) (0.2125,0.55911) (0.225,0.58322) (0.2375,0.60556) (0.25,0.6256) (0.2625,0.64285) (0.275,0.65686) (0.2875,0.66725) (0.3,0.67371) (0.3125,0.67607) (0.325,0.67427) (0.3375,0.6684) (0.35,0.65859) (0.3625,0.64508) (0.375,0.62811) (0.3875,0.60797) (0.4,0.58496) (0.4125,0.5594) (0.425,0.53158) (0.4375,0.50183) (0.45,0.47045) (0.4625,0.43775) (0.475,0.40403) (0.4875,0.3696) (0.5,0.33473) (0.5125,0.29973) (0.525,0.26487) (0.5375,0.23043) (0.55,0.19671) (0.5625,0.164) (0.575,0.13261) (0.5875,0.10288) (0.6,0.075137) (0.6125,0.049749) (0.625,0.027083) (0.6375,0.0075115) (0.65,-0.0085968) (0.6625,-0.020893) (0.675,-0.029065) (0.6875,-0.032865) (0.7,-0.032138) (0.7125,-0.026867) (0.725,-0.017196) (0.7375,-0.0034041) (0.75,0.014132) (0.7625,0.034972) (0.775,0.058639) (0.7875,0.084635) (0.8,0.11246) (0.8125,0.14161) (0.825,0.17161) (0.8375,0.202) (0.85,0.23236) (0.8625,0.26228) (0.875,0.29141) (0.8875,0.31943) (0.9,0.34611) (0.9125,0.37123) (0.925,0.39469) (0.9375,0.41643) (0.95,0.43644) (0.9625,0.45475) (0.975,0.4714) (0.9875,0.48647) (1,0.50004)};\addlegendentry{prédiction $\mu$}
\addplot[only marks,mark=*,mark size=2pt,cink] coordinates {(-1,1.1676) (0.1,0.34879) (0.3,0.67431) (0.5,0.33467) (0.7,-0.032945) (0.9,0.34634)};\addlegendentry{simulations}
\addlegendimage{area legend,fill=cblue!15,draw=none}\addlegendentry{$\mu\pm2s$}
\nextgroupplot[height=3.6cm,xlabel={$x$},ylabel={EI},ymin=0]
\addplot[corange,very thick,fill=corange!15] coordinates {(-1,0) (-0.9875,4.9407e-324) (-0.975,2.7825e-97) (-0.9625,6.1889e-47) (-0.95,1.8507e-28) (-0.9375,1.1589e-19) (-0.925,9.0799e-15) (-0.9125,9.509e-12) (-0.9,9.6534e-10) (-0.8875,2.4749e-08) (-0.875,2.6666e-07) (-0.8625,1.6166e-06) (-0.85,6.5826e-06) (-0.8375,2.0155e-05) (-0.825,5.0012e-05) (-0.8125,0.00010586) (-0.8,0.00019818) (-0.7875,0.00033685) (-0.775,0.00053001) (-0.7625,0.00078329) (-0.75,0.0010996) (-0.7375,0.001479) (-0.725,0.0019194) (-0.7125,0.0024165) (-0.7,0.0029646) (-0.6875,0.0035571) (-0.675,0.0041868) (-0.6625,0.0048464) (-0.65,0.0055286) (-0.6375,0.0062266) (-0.625,0.0069339) (-0.6125,0.0076449) (-0.6,0.0083545) (-0.5875,0.0090583) (-0.575,0.0097526) (-0.5625,0.010435) (-0.55,0.011102) (-0.5375,0.011753) (-0.525,0.012386) (-0.5125,0.013) (-0.5,0.013596) (-0.4875,0.014172) (-0.475,0.01473) (-0.4625,0.01527) (-0.45,0.015792) (-0.4375,0.016296) (-0.425,0.016785) (-0.4125,0.017258) (-0.4,0.017716) (-0.3875,0.01816) (-0.375,0.01859) (-0.3625,0.019005) (-0.35,0.019407) (-0.3375,0.019793) (-0.325,0.020164) (-0.3125,0.020517) (-0.3,0.020851) (-0.2875,0.021161) (-0.275,0.021445) (-0.2625,0.021697) (-0.25,0.021913) (-0.2375,0.022083) (-0.225,0.022201) (-0.2125,0.022257) (-0.2,0.022239) (-0.1875,0.022134) (-0.175,0.021928) (-0.1625,0.021606) (-0.15,0.021151) (-0.1375,0.020545) (-0.125,0.019771) (-0.1125,0.018811) (-0.1,0.017653) (-0.0875,0.016286) (-0.075,0.014709) (-0.0625,0.012932) (-0.05,0.010981) (-0.0375,0.0089056) (-0.025,0.0067865) (-0.0125,0.0047384) (0,0.0029102) (0.0125,0.0014658) (0.025,0.00053123) (0.0375,0.00010581) (0.05,6.1267e-06) (0.0625,1.6632e-08) (0.075,1.8944e-15) (0.0875,5.6547e-46) (0.1,9.6857e-202) (0.1125,1.4024e-60) (0.125,6.0882e-25) (0.1375,6.7459e-16) (0.15,1.8143e-12) (0.1625,5.6279e-11) (0.175,2.0319e-10) (0.1875,1.8176e-10) (0.2,4.7796e-11) (0.2125,3.0384e-12) (0.225,2.4811e-14) (0.2375,6.0236e-18) (0.25,1.2849e-24) (0.2625,1.2475e-38) (0.275,1.6657e-75) (0.2875,3.0802e-218) (0.3,0) (0.3125,1.1522e-226) (0.325,3.7575e-81) (0.3375,8.0495e-43) (0.35,4.7543e-28) (0.3625,5.8348e-21) (0.375,4.0228e-17) (0.3875,6.0933e-15) (0.4,9.2582e-14) (0.4125,2.65e-13) (0.425,1.6357e-13) (0.4375,1.5727e-14) (0.45,7.7327e-17) (0.4625,7.7879e-22) (0.475,2.398e-34) (0.4875,8.1716e-78) (0.5,2.8601e-187) (0.5125,1.1494e-54) (0.525,3.806e-18) (0.5375,2.4607e-09) (0.55,3.9277e-06) (0.5625,0.00012331) (0.575,0.00081289) (0.5875,0.0025484) (0.6,0.005328) (0.6125,0.0087022) (0.625,0.012027) (0.6375,0.014652) (0.65,0.016003) (0.6625,0.015624) (0.675,0.013215) (0.6875,0.0088286) (0.7,0.0046828) (0.7125,0.0062182) (0.725,0.008678) (0.7375,0.0099115) (0.75,0.0098572) (0.7625,0.0087625) (0.775,0.0069754) (0.7875,0.0048907) (0.8,0.0029064) (0.8125,0.0013601) (0.825,0.00043488) (0.8375,7.0752e-05) (0.85,2.9323e-06) (0.8625,4.39e-09) (0.875,1.5126e-16) (0.8875,3.1778e-48) (0.9,1.9054e-198) (0.9125,1.7517e-54) (0.925,1.5594e-20) (0.9375,6.0629e-12) (0.95,1.6555e-08) (0.9625,9.5406e-07) (0.975,1.063e-05) (0.9875,5.1646e-05) (1,0.00015664)} \closedcycle;
\end{groupplot}
\end{tikzpicture}

\caption{Krigeage sur un exemple à une variable, $f(x)=0{,}6\sin(6x)+x^2$. En haut : la fonction (tirets), six simulations (points), la prédiction $\mu$ (trait bleu) et la bande $\mu\pm2s$. L'incertitude s'annule presque aux points simulés et devient grande dans l'intervalle de gauche, où rien n'a été simulé. En bas : l'amélioration espérée du chapitre~\ref{ch:acquisition}.}\label{fig:gp1d}
\end{figure}

\section{Les hyperparamètres}
Le modèle dépend de $m$, $\sigma^2$, des $\theta_i$ et du bruit $\sigma_n^2$. On les choisit par \textbf{maximum de vraisemblance} : on retient les valeurs pour lesquelles les simulations observées sont les plus probables,
\[
\log p(\vect y\mid X)=-\tfrac12(\vect y-m\vect 1)^{\!\top}K^{-1}(\vect y-m\vect 1)-\tfrac12\log\det K-\tfrac n2\log 2\pi .
\]
Le premier terme récompense l'accord avec les données, le deuxième pénalise les modèles trop souples. Pour $m$ et $\sigma^2$, le maximum a une forme explicite ; pour les $\theta_i$ et le bruit, il faut une optimisation numérique, refaite à chaque itération de la boucle.

\begin{piege}
Chaque entraînement demande plusieurs factorisations de $K$, en $O(n^3)$. Avec 250 profils et 4 variables, c'est environ 1~s par modèle sur un ordinateur portable, contre quelques millisecondes pour un modèle quadratique. Sur la campagne complète, les variantes à krigeage coûtent 25 à 30 fois plus de calcul que les variantes quadratiques (chapitre~\ref{ch:resultats}).
\end{piege}

\section{Pourquoi estimer un bruit}
XFOIL est déterministe : relancer un profil donne le même résultat. Pourtant, on estime un bruit $\sigma_n^2$. La raison est que la traînée n'est pas une fonction lisse des paramètres : la position de la transition laminaire-turbulent bascule d'un profil à l'autre, et $C_{L,\max}$ est estimé sur une rampe d'incidence discrète. Vues à l'échelle de la grille, ces irrégularités se comportent comme un bruit. Sans terme de bruit, le krigeage chercherait à interpoler chaque bascule, avec des $\theta_i$ très grands et une prédiction dégradée entre les points.

\section{Le modèle quadratique vu comme un krigeage}
Le modèle quadratique peut lui aussi fournir une incertitude. Notons $\varphi(x)$ le vecteur des monômes de degré au plus 2 (15 monômes pour 4 variables). On suppose $y=\varphi(x)^{\!\top}\vect c+\varepsilon$, avec un a priori $\vect c\sim\mathcal N(0,\lambda^{-1}I)$ sur les coefficients et un bruit de variance $\alpha^{-1}$. C'est la \textbf{régression bayésienne}. La prédiction en $x$ est normale, de moyenne $\varphi(x)^{\!\top}\hat{\vect c}$ et de variance
\[
s^2(x)=\alpha^{-1}+\varphi(x)^{\!\top}\Sigma\,\varphi(x),\qquad \Sigma=\big(\lambda I+\alpha\,\Phi^{\!\top}\Phi\big)^{-1},
\]
où $\Phi$ est la matrice des monômes aux points simulés. Les deux précisions $\alpha$ et $\lambda$ sont réglées par maximum de vraisemblance, comme pour le krigeage.

\begin{resultat}
La régression quadratique bayésienne \emph{est} un processus gaussien, de noyau $k(x,x')=\lambda^{-1}\varphi(x)^{\!\top}\varphi(x')$. La différence avec le krigeage n'est donc pas « bayésien ou non », mais la famille de fonctions : 15 coefficients d'un côté, une famille bien plus large de l'autre.
\end{resultat}

Cette remarque rend l'ablation possible : la case « modèle quadratique + amélioration espérée » du chapitre~\ref{ch:question} utilise cette variance.

\begin{codebox}[ : \code{optim/bayes.py}]
\begin{lstlisting}
class GPModel:      # krigeage SMT : Matern 5/2, moyenne constante, bruit estime
    def __init__(self, Z, y):
        self.gp = KRG(theta0=[0.5] * d, corr="matern52", poly="constant",
                      eval_noise=True, noise0=[1e-3], n_start=1, hyper_opt="Cobyla")
        self.gp.set_training_values(Z, ((y - self.mu) / self.sd).reshape(-1, 1))
        self.gp.train()
    def predict(self, Z):                    # -> (moyenne, ecart-type)
        m = self.gp.predict_values(Z)[:, 0] * self.sd + self.mu
        v = np.maximum(self.gp.predict_variances(Z)[:, 0], 0.0)
        return m, np.sqrt(v) * self.sd

class QuadModel:    # regression quadratique bayesienne (sklearn BayesianRidge)
    def predict(self, Z):
        return self.br.predict(quadratic_features(Z)[:, 1:], return_std=True)
\end{lstlisting}
Les deux classes ont la même interface : la boucle ne sait pas lequel des deux modèles elle utilise. Les variables sont normalisées dans $[-1,1]$, comme dans la boucle QUBO. Avec moins de 4 points ou une sortie constante, on se replie sur un modèle constant.
\end{codebox}

\exercices
\begin{exercice}[un seul point simulé]
On a simulé un seul profil $x_1$, de valeur $y_1$, sans bruit, avec $m=0$. Écrire $\mu(x)$ et $s^2(x)$ en fonction de $k$. Que valent-ils en $x=x_1$ ? Et très loin de $x_1$, quand $k(x_1,x)\to0$ ?
\end{exercice}
\begin{exercice}
Dans l'équation~\eqref{eq:krigeage}, pourquoi $s^2(x)$ ne dépend-il pas de $\vect y$ ? Est-ce encore vrai une fois les hyperparamètres réglés par maximum de vraisemblance ?
\end{exercice}
\begin{exercice}[noyau polynomial]
Montrer que si $\vect c\sim\mathcal N(0,\lambda^{-1}I)$ et $f(x)=\varphi(x)^{\!\top}\vect c$, alors $\mathrm{Cov}(f(x),f(x'))=\lambda^{-1}\varphi(x)^{\!\top}\varphi(x')$.
\end{exercice}

\begin{pistes}
\begin{itemize}
\item Utiliser les noyaux de SMT prévus pour les variables entières (relaxation continue, distance de Gower) au lieu de traiter les indices comme des réels.
\item Krigeage avec réduction de dimension (KPLS) si la grille grandit en nombre de variables.
\item Mettre à jour le krigeage sans le réentraîner entièrement à chaque itération (hyperparamètres gelés quelques itérations), pour réduire le coût de calcul.
\end{itemize}
\end{pistes}

\chapter{L'amélioration espérée : choisir le prochain profil}\label{ch:acquisition}
\iterbanner{V3}

\section{Exploiter ou explorer}
Une fois le modèle ajusté, il faut choisir le profil à simuler. Deux stratégies extrêmes :
\begin{itemize}
\item \textbf{exploiter} : simuler le profil dont la prédiction $\mu(x)$ est la plus basse. C'est ce que fait la boucle QUBO. Risque : si le modèle se trompe sur une région jamais visitée, on ne le saura jamais.
\item \textbf{explorer} : simuler le profil dont l'incertitude $s(x)$ est la plus grande. On améliore le modèle partout, y compris là où il n'y a rien à gagner.
\end{itemize}
Un \emph{critère d'acquisition} combine les deux en un seul nombre par profil. On simule le profil qui le maximise.

\section{L'amélioration espérée}
On minimise une fonction $f$. Soit $f_{\min}$ la meilleure valeur déjà simulée. Pour un profil $x$ non simulé, le modèle dit $f(x)\sim\mathcal N(\mu,s^2)$. L'\emph{amélioration} qu'apporterait sa simulation est $\max(f_{\min}-f(x),\,0)$ : ce qu'on gagne s'il est meilleur, zéro sinon.

\begin{definition}[ : amélioration espérée]
\[
\mathrm{EI}(x)=\mathbb E\big[\max(f_{\min}-f(x),\,0)\big]
=(f_{\min}-\mu)\,\Phi(u)+s\,\varphi_{\mathcal N}(u),\qquad u=\frac{f_{\min}-\mu}{s},
\]
où $\Phi$ et $\varphi_{\mathcal N}$ sont la fonction de répartition et la densité de la loi normale centrée réduite.
\end{definition}

\paragraph{D'où vient la formule.} On écrit $f(x)=\mu+s\,\xi$ avec $\xi\sim\mathcal N(0,1)$. L'amélioration est non nulle quand $\xi<u$, et vaut alors $s(u-\xi)$ :
\[
\mathrm{EI}=s\int_{-\infty}^{u}(u-\xi)\,\varphi_{\mathcal N}(\xi)\,\dd\xi
=s\,u\,\Phi(u)-s\int_{-\infty}^{u}\xi\,\varphi_{\mathcal N}(\xi)\,\dd\xi
=s\,u\,\Phi(u)+s\,\varphi_{\mathcal N}(u),
\]
car $\varphi_{\mathcal N}'(\xi)=-\xi\,\varphi_{\mathcal N}(\xi)$. On retrouve la formule avec $s\,u=f_{\min}-\mu$.

\begin{resultat}
\begin{itemize}
\item Le premier terme, $(f_{\min}-\mu)\Phi(u)$, est grand quand la prédiction est \emph{meilleure} que le record : c'est l'exploitation.
\item Le second, $s\,\varphi_{\mathcal N}(u)$, est grand quand l'incertitude est grande : c'est l'exploration.
\item $\mathrm{EI}\ge0$, et $\mathrm{EI}=0$ en un point déjà simulé sans bruit ($s=0$, $\mu\ge f_{\min}$) : on ne resimule jamais.
\item Quand $s\to0$, $\mathrm{EI}\to\max(f_{\min}-\mu,0)$ : sans incertitude, l'amélioration espérée redonne la stratégie « minimum de la prédiction ».
\end{itemize}
\end{resultat}

La figure~\ref{fig:gp1d} (en bas) montre le critère sur l'exemple à une variable. Il a deux bosses. Celle de droite entoure le minimum de la prédiction ($x\approx0{,}69$) : c'est l'exploitation. Celle de gauche, la plus haute, est au milieu de l'intervalle sans simulation ($x\approx-0{,}21$) : c'est l'exploration. La stratégie « minimum de la prédiction » simulerait $x\approx0{,}69$, près d'un minimum local déjà bien connu. L'amélioration espérée simule $x\approx-0{,}21$, à côté du vrai minimum global ($x\approx-0{,}24$), que le modèle ne voit pas encore.

\section{L'algorithme EGO}
L'optimisation bayésienne à un objectif, sous sa forme classique (\emph{Efficient Global Optimization}, Jones, Schonlau et Welch, 1998), est la boucle :
\begin{enumerate}
\item simuler quelques points initiaux ;
\item ajuster le krigeage ;
\item simuler le point qui maximise l'amélioration espérée ;
\item recommencer en 2 jusqu'à épuisement du budget.
\end{enumerate}

\begin{projet}
En variables continues, l'étape 3 est elle-même une optimisation difficile, car l'amélioration espérée a de nombreux maxima locaux (on la traite par descentes multiples). Notre espace de conception est une grille de 1\,024 profils : on calcule simplement le critère sur \emph{tous} les profils non simulés et on prend le meilleur. L'étape 3 coûte une prédiction du modèle sur la grille, soit quelques centièmes de seconde.
\end{projet}

\begin{codebox}[ : \code{optim/bayes.py}]
\begin{lstlisting}
def expected_improvement(mean, std, best):
    """E[max(best - Y, 0)], Y ~ N(mean, std^2) : amelioration esperee pour une minimisation."""
    std = np.maximum(std, 1e-12)
    u = (best - mean) / std
    return (best - mean) * norm.cdf(u) + std * norm.pdf(u)
\end{lstlisting}
Vérification : pour $\mu=f_{\min}$ et $s=1$, on doit trouver $\varphi_{\mathcal N}(0)=1/\sqrt{2\pi}\approx0{,}3989$.
\end{codebox}

\section{D'autres critères}
L'amélioration espérée n'est pas le seul critère. Deux autres sont courants :
\begin{itemize}
\item la \textbf{borne de confiance inférieure}, $\mu(x)-\kappa\,s(x)$, à minimiser : simple, mais le poids $\kappa$ de l'exploration est à régler ;
\item la \textbf{probabilité d'amélioration}, $\Phi(u)$ : elle ignore l'ampleur du gain et tend à trop exploiter.
\end{itemize}
On s'en tient à l'amélioration espérée, qui n'a pas de paramètre à régler.

\exercices
\begin{exercice}
Deux profils non simulés ont pour prédictions $(\mu,s)=(0{,}50;\ 0{,}02)$ et $(0{,}60;\ 0{,}20)$, et le record est $f_{\min}=0{,}52$. Calculer l'amélioration espérée de chacun. Lequel simule-t-on ? Lequel aurait choisi la stratégie « minimum de la prédiction » ?
\end{exercice}
\begin{exercice}
Montrer que l'amélioration espérée est croissante en $s$ à $\mu$ fixé, et décroissante en $\mu$ à $s$ fixé. (Indication : $\partial\mathrm{EI}/\partial s=\varphi_{\mathcal N}(u)$ et $\partial\mathrm{EI}/\partial\mu=-\Phi(u)$.)
\end{exercice}

\begin{pistes}
\begin{itemize}
\item Essayer les critères de l'ONERA (WB2, WB2S), qui ajoutent à l'amélioration espérée un terme proportionnel à la prédiction pour faciliter sa maximisation en variables continues.
\item Sur la grande grille du lot L2, mesurer à partir de quelle taille l'énumération du critère devient le poste de calcul dominant : c'est là que l'hypothèse H2 situe l'intérêt d'une sélection par QUBO.
\end{itemize}
\end{pistes}

\chapter{Contraintes et deux objectifs}\label{ch:contraintes}
\iterbanner{V3}

Le chapitre précédent traite un objectif sans contrainte. Notre problème a deux objectifs (traînée et SER) et deux contraintes (épaisseur et manœuvre). Ce chapitre décrit comment on étend la méthode, et quels choix on a retenus.

\section{La contrainte connue : l'épaisseur}
$t\ge8$~\% ne dépend que d'une variable de conception : on la connaît sans rien simuler. La boucle QUBO l'encode par une pénalité et un qubit auxiliaire (réduction de Rosenberg). Sur une grille énumérée, il suffit de \textbf{retirer} de la liste des candidats les profils d'indice d'épaisseur nul. Le résultat est le même, sans qubit.

\section{La contrainte simulée : la manœuvre}
$C_{L,\max}\ge1{,}2$ n'est connue qu'après simulation. On la modélise donc comme un objectif : un modèle de $C_{L,\max}$, de la même famille que ceux des objectifs, donne pour chaque profil une prédiction $\mu_c(x)$ et un écart-type $s_c(x)$. Reste à décider comment s'en servir. Quatre traitements sont codés.

\begin{center}
\renewcommand{\arraystretch}{1.3}
\begin{tabularx}{\linewidth}{>{\raggedright\arraybackslash}p{3.1cm}X>{\raggedright\arraybackslash}p{3.6cm}}
\toprule
Traitement & Règle & Nom dans le code \\
\midrule
en dur sur la prédiction & on ne garde que les profils tels que $\mu_c(x)\ge1{,}2$ (c'est le principe de l'algorithme SEGO) & \code{constraint="mean"} \\
probabilité de faisabilité & on multiplie l'amélioration espérée par $\mathrm{PoF}(x)$ & \code{"pof"}, sélection \code{"ei"} \\
pénalité probabiliste & on ajoute $1-\mathrm{PoF}(x)$ à la prédiction de l'objectif & \code{"pof"}, sélection \code{"mean"} \\
pénalité de la boucle QUBO & modèle \emph{affine} de $C_{L,\max}$ et pénalité $75\,\max(1{,}2-\widehat{C_{L,\max}},\,0)^2$ & \code{"lin"} \\
\bottomrule
\end{tabularx}
\end{center}

\begin{definition}[ : probabilité de faisabilité]
Si le modèle dit $C_{L,\max}(x)\sim\mathcal N(\mu_c,s_c^2)$, la probabilité que le profil tienne la contrainte est
\[
\mathrm{PoF}(x)=\mathbb P\big(C_{L,\max}(x)\ge1{,}2\big)=\Phi\!\left(\frac{\mu_c(x)-1{,}2}{s_c(x)}\right).
\]
\end{definition}

Le critère $\mathrm{EI}(x)\times\mathrm{PoF}(x)$ est l'amélioration espérée \emph{sous contrainte}, si l'on suppose l'objectif et la contrainte indépendants : on ne gagne quelque chose que si le profil est à la fois meilleur et faisable.

\paragraph{La pénalité de la boucle QUBO.} Dans \code{qubo\_loop.py}, la contrainte devient $P_c\,(\widehat{C_{L,\max}}-1{,}2-s)^2$ avec une variable d'écart $s\ge0$ sur 3 qubits et $P_c=3\times25=75$. Minimisée sur $s$, cette expression vaut à peu près $75\,\max(1{,}2-\widehat{C_{L,\max}},0)^2$ : rien si le profil est prédit faisable, un coût quadratique sinon. Le traitement \code{"lin"} applique directement cette formule. Il sert au \textbf{contrôle} de l'ablation : même modèle affine, même poids, sans encodage binaire.

\begin{piege}
Ces traitements ne sont pas des variantes de détail. Sur notre problème, le front exact est collé à la contrainte de manœuvre, et le choix du traitement décide du succès ou de l'échec de la méthode, bien plus que le choix du modèle. C'est l'objet du chapitre~\ref{ch:diagnostic}.
\end{piege}

\section{Deux objectifs : la scalarisation}
L'amélioration espérée suppose un seul objectif. On se ramène à ce cas par une \emph{scalarisation} : pour chaque poids $w$, on fabrique un objectif unique à partir des deux, et on parcourt plusieurs poids pour couvrir le front. Les objectifs sont d'abord normalisés par leur étendue sur les profils faisables déjà simulés : $\tilde f_1$ pour la traînée, $\tilde f_2$ pour la SER.

\paragraph{Somme pondérée.} $g_w=w\,\tilde f_1+(1-w)\,\tilde f_2$, comme dans la boucle QUBO. Avantage : si les deux objectifs ont chacun leur modèle, de lois $\mathcal N(\mu_1,s_1^2)$ et $\mathcal N(\mu_2,s_2^2)$ supposées indépendantes, alors $g_w$ est lui aussi normal :
\[
g_w(x)\sim\mathcal N\Big(w\mu_1+(1-w)\mu_2,\ \ w^2s_1^2+(1-w)^2s_2^2\Big).
\]
Deux modèles suffisent pour tous les poids, et l'amélioration espérée de $g_w$ se calcule par la formule du chapitre~\ref{ch:acquisition}. Limite connue (cours d'optimisation, chapitre sur la dominance) : une somme pondérée n'atteint que les points de l'enveloppe convexe du front.

\paragraph{Tchebychev augmenté (ParEGO).} L'algorithme ParEGO (Knowles, 2006) utilise
\[
g_w=\max\big(w\,\tilde f_1,\ (1-w)\,\tilde f_2\big)+\rho\,\big(w\,\tilde f_1+(1-w)\,\tilde f_2\big),\qquad\rho=0{,}05,
\]
qui atteint aussi les points situés dans les creux du front. Contrepartie : le maximum de deux lois normales n'est pas normal. ParEGO ajuste donc un krigeage directement sur les valeurs scalarisées, \emph{un par poids} : cinq modèles par itération au lieu de deux.

\begin{projet}
On compare les deux : la somme pondérée pour l'ablation (même scalarisation que la boucle QUBO, donc une seule chose change à la fois), et ParEGO comme variante.
\end{projet}

\section{Un lot de cinq profils par itération}
L'algorithme EGO simule un profil par itération. La boucle QUBO en simule cinq, un par poids. Pour comparer à structure égale, on fait de même : pour chaque poids, on classe les candidats selon le critère et on prend le meilleur qui n'est ni déjà simulé, ni déjà retenu pour un autre poids dans la même itération. Le modèle n'est réajusté qu'une fois les cinq profils simulés.

\begin{resultat}
Les variantes comparées ne diffèrent de la boucle QUBO que par le modèle, la sélection, et le traitement de la contrainte de manœuvre. Points initiaux, poids, taille du lot, compteur et budget sont identiques.
\end{resultat}

\section{Ce qu'on n'a pas fait : l'amélioration de l'hypervolume}
Une autre voie évite la scalarisation : choisir le profil qui maximise l'\emph{amélioration espérée de l'hypervolume} du front (EHVI). C'est le critère de référence en optimisation bayésienne multi-objectif ; il se calcule exactement pour deux objectifs. Il n'est pas codé ici : la somme pondérée et ParEGO suffisent à répondre à H1, et gardent la structure « un profil par poids » de la boucle QUBO.

\exercices
\begin{exercice}
Un profil a une amélioration espérée de $0{,}04$, et le modèle de contrainte prédit $C_{L,\max}=1{,}15\pm0{,}10$. Calculer sa probabilité de faisabilité et son critère $\mathrm{EI}\times\mathrm{PoF}$. Serait-il candidat avec le traitement en dur ?
\end{exercice}
\begin{exercice}
Pour $w=0{,}5$, $\mu_1=0{,}3$, $s_1=0{,}1$, $\mu_2=0{,}5$, $s_2=0{,}2$, donner la loi de la somme pondérée. Pourquoi l'hypothèse d'indépendance des deux objectifs est-elle discutable pour la traînée et la SER d'un même profil ?
\end{exercice}
\begin{exercice}
Avec la pénalité de la boucle QUBO, quel est le coût d'un profil prédit à $\widehat{C_{L,\max}}=1{,}10$ ? Le comparer à l'étendue des objectifs normalisés (de l'ordre de 1). Même question pour $1{,}15$.
\end{exercice}

\begin{pistes}
\begin{itemize}
\item Coder l'amélioration espérée de l'hypervolume et la comparer à ParEGO.
\item Modéliser la corrélation entre traînée, SER et $C_{L,\max}$ (krigeage multi-sorties) au lieu de les supposer indépendants.
\item Sélectionner le lot de cinq profils en tenant compte de leur redondance (critère « q-EI ») au lieu d'un profil par poids.
\end{itemize}
\end{pistes}

\chapter{Le code et sa validation}\label{ch:code}
\iterbanner{V3}

\section{Trois fichiers}
\begin{center}
\renewcommand{\arraystretch}{1.2}
\begin{tabularx}{\linewidth}{>{\raggedright\arraybackslash}p{4.4cm}X}
\toprule
Fichier & Rôle \\
\midrule
\code{optim/bayes.py} & les deux modèles (\code{GPModel}, \code{QuadModel}), l'amélioration espérée, et la boucle \code{bayes\_loop} \\
\code{optim/experiences\_l1.py} & la comparaison : 8 méthodes $\times$ 2 scénarios $\times$ 20 graines $\to$ \code{resultats/pipeline\_L1\_V3.json} \\
\code{optim/diagnostic\_l1.py} & les diagnostics du chapitre~\ref{ch:diagnostic} $\to$ \code{resultats/diagnostic\_L1\_V3.json} \\
\bottomrule
\end{tabularx}
\end{center}
Aucun fichier de la V2 n'est modifié : \code{bayes\_loop} réutilise \code{CountedProblem}, la normalisation \code{z\_of} et le protocole de \code{experiences\_v2.py}.

\section{La boucle}
\code{bayes\_loop} a la même signature et la même structure que \code{qubo\_loop}. Trois arguments choisissent la variante :
\begin{center}
\begin{tabular}{lll}
\toprule
Argument & Valeurs & Chapitre \\
\midrule
\code{model} & \code{"gp"} (krigeage), \code{"quad"} (quadratique bayésien) & \ref{ch:krigeage} \\
\code{selection} & \code{"ei"} (amélioration espérée), \code{"mean"} (minimum de la prédiction) & \ref{ch:acquisition} \\
\code{scalar} & \code{"sum"} (somme pondérée), \code{"tcheby"} (ParEGO) & \ref{ch:contraintes} \\
\code{constraint} & \code{"pof"}, \code{"mean"}, \code{"lin"} & \ref{ch:contraintes} \\
\bottomrule
\end{tabular}
\end{center}

\begin{codebox}[ : le cœur d'une itération (\code{optim/bayes.py}, simplifié)]
\begin{lstlisting}
m1, s1 = Model(Z[fit], Fn[fit, 0]).predict(Zall)      # trainee normalisee, sur toute la grille
m2, s2 = Model(Z[fit], Fn[fit, 1]).predict(Zall)      # SER normalisee
mc, sc = Model(Z[vl], cl[vl]).predict(Zall)           # C_L,max
pof = norm.cdf((mc - cl_min) / sc)                    # probabilite de faisabilite
for w in weights:                                     # 5 poids -> 5 profils
    y_obs = w * Fn[:, 0] + (1 - w) * Fn[:, 1]
    mean = w * m1 + (1 - w) * m2
    std = np.sqrt((w * s1) ** 2 + ((1 - w) * s2) ** 2)
    if selection == "ei":
        score = expected_improvement(mean, std, y_obs[ok].min()) * pof
    else:
        score = -mean - (1.0 - pof)
    # meilleur profil non simule, d'epaisseur permise, pas deja retenu pour un autre poids
\end{lstlisting}
\end{codebox}

Trois détails comptent pour l'équité de la comparaison :
\begin{itemize}
\item le générateur aléatoire est initialisé et consommé comme dans \code{qubo\_loop} : pour une graine donnée, les 10 profils initiaux sont identiques ;
\item les objectifs sont ajustés sur les profils \emph{faisables} déjà simulés, et $C_{L,\max}$ sur les profils dont la croisière a convergé, comme dans la boucle QUBO ;
\item tant qu'il y a moins de 4 profils faisables, les deux boucles explorent au hasard, de la même façon.
\end{itemize}

\section{Validation}
\paragraph{1. L'amélioration espérée.} Pour $\mu=f_{\min}$ et $s=1$, le code rend $0{,}3989$, soit $1/\sqrt{2\pi}$.

\paragraph{2. Les modèles, sur un cas jouet.} Fonction lisse de 4 variables, $f(z)=\sin(2z_1)+z_2^2-\tfrac12z_3z_4$, 60 points d'entraînement bruités, 500 points de test. L'écart-type de $f$ sur le test est $0{,}88$.
\begin{center}
\begin{tabular}{lcc}
\toprule
 & erreur quadratique moyenne & erreurs dans $\pm2s$ \\
\midrule
krigeage & 0,039 & 99,8~\% \\
modèle quadratique & 0,18 & 92~\% \\
\bottomrule
\end{tabular}
\end{center}
Le krigeage est près de cinq fois plus précis sur cette fonction, qui n'est pas quadratique (le sinus). Pour les deux modèles, l'incertitude annoncée est cohérente avec l'erreur réelle (on attend environ 95~\%).

\paragraph{3. La boucle QUBO de la V2 est reproduite.} \code{experiences\_l1.py} relance la boucle QUBO à solveur exact sur les 20 graines. On retrouve exactement les médianes de la fiche projet : 75 simulations dans S1 et 142,5 dans S2. La chaîne tourne donc à l'identique sur une autre machine (Windows, sans XFOIL, à partir de la table \code{base\_V2.csv}).

\paragraph{4. Le contrôle.} La variante « modèle quadratique + minimum de la prédiction + pénalité de la boucle QUBO » doit se comporter comme la boucle QUBO. C'est le cas dans S1 (médiane de 75 simulations pour les deux). Dans S2 elle fait mieux (95 contre 142,5) : le chapitre~\ref{ch:resultats} y revient.

\begin{piege}
La première version du contrôle imposait la contrainte de manœuvre en dur sur la prédiction. Elle ne retrouvait \emph{jamais} le front, alors que la boucle QUBO y arrive toujours. Ce n'était pas une erreur de code mais une vraie différence de méthode ; la comprendre a donné le résultat principal du lot (chapitre~\ref{ch:diagnostic}). Un contrôle qui échoue est une information : il ne faut pas le retirer du tableau.
\end{piege}

\section{Lancer les calculs}
\begin{lstlisting}
cd Code
python -m optim.experiences_l1 rapide S1    # essai : 4 graines, un scenario (4 min)
python -m optim.experiences_l1              # campagne : 20 graines, S1 et S2 (19 min sur 10 processus)
python -m optim.diagnostic_l1               # diagnostics (5 min)
cd ../Cours/5_Optimisation_Bayesienne
python tools/make_figures.py                # figures et chiffres du cours
\end{lstlisting}
Dépendances ajoutées par ce lot : \code{smt} (krigeage) et \code{scikit-learn} (régression bayésienne).

\exercices
\begin{exercice}
Pourquoi faut-il que les 10 profils initiaux soient les mêmes pour toutes les méthodes à graine donnée ? Que gagne-t-on pour l'analyse statistique ?
\end{exercice}
\begin{exercice}
Dans l'extrait de code, les objectifs normalisés \code{Fn} sont de l'ordre de 0 à 1. Quelle est alors la signification de la pénalité \code{(1.0 - pof)} ajoutée pour la sélection \code{"mean"} ?
\end{exercice}

\begin{pistes}
\begin{itemize}
\item Ajouter des tests automatiques (amélioration espérée, cas jouet, reproduction des médianes V2) à \code{verifier\_installation.py}.
\item \iterbadge{V3}~\textbf{En partie} : \code{smt} et \code{scikit-learn} sont dans \code{requirements.txt}. Reste \code{installer\_outils.sh}, qui n'installe encore que \code{pymoo} dans \code{.venv}.
\end{itemize}
\end{pistes}

\chapter{Résultats : la boucle QUBO face à l'optimisation bayésienne}\label{ch:resultats}
\iterbanner{V3}

\section{Les méthodes comparées}
\begin{center}
\renewcommand{\arraystretch}{1.2}
\begin{tabularx}{\linewidth}{>{\raggedright\arraybackslash}p{3.9cm}llX}
\toprule
Méthode & Modèle & Sélection & Contrainte de manœuvre \\
\midrule
boucle QUBO (V2) & quadratique & minimum (QUBO exact) & modèle affine, pénalité à qubits d'écart \\
quadratique + minimum & quadratique & minimum & modèle affine, même pénalité, sans qubits \\
krigeage + minimum & krigeage & minimum & idem \\
quadratique + EI & quadratique & amélioration espérée & probabilité de faisabilité \\
krigeage + EI & krigeage & amélioration espérée & probabilité de faisabilité \\
ParEGO & krigeage par poids & amélioration espérée & probabilité de faisabilité \\
\midrule
quadratique + minimum, contrainte stricte & quadratique & minimum & pénalité probabiliste $1-\mathrm{PoF}$ \\
krigeage + EI, contrainte en dur & krigeage & amélioration espérée & $\mu_c\ge1{,}2$ exigé \\
\bottomrule
\end{tabularx}
\end{center}
Les quatre lignes « quadratique/krigeage $\times$ minimum/EI » forment l'ablation du chapitre~\ref{ch:question}. Les deux dernières lignes ne changent que le traitement de la contrainte.

\section{Retrouver tout le front exact}
\begin{figure}[htbp]
\centering
% figure générée par tools/make_figures.py : l1_sims
\begin{tikzpicture}
\begin{groupplot}[group style={group size=2 by 1,horizontal sep=0.5cm},width=0.4\linewidth,height=6.4cm,
  xlabel={simulations pour le front complet},xmin=0,xmax=250,ymin=0.4,ymax=8.6,
  title style={font=\small},tick label style={font=\scriptsize},label style={font=\small},
  yticklabel style={font=\scriptsize,align=right},xmajorgrids,grid style={cmuted!25}]
\nextgroupplot[title={S1 : approche de face},ytick={8,7,6,5,4,3,2,1},yticklabels={{boucle QUBO (V2)},{quadratique + minimum},{krigeage + minimum},{quadratique + EI},{krigeage + EI},{ParEGO},{quadratique + minimum, contrainte stricte},{krigeage + EI, contrainte en dur}}]
\addplot[cink,line width=2.2pt,forget plot] coordinates {(60,8) (96.25,8)};
\addplot[only marks,mark=*,mark size=2.4pt,cink,mark options={fill=white,line width=1pt},forget plot] coordinates {(75,8)};
\addplot[cblue,line width=2.2pt,forget plot] coordinates {(70,7) (111.2,7)};
\addplot[only marks,mark=*,mark size=2.4pt,cblue,mark options={fill=white,line width=1pt},forget plot] coordinates {(75,7)};
\addplot[cblue,line width=2.2pt,forget plot] coordinates {(75,6) (105,6)};
\addplot[only marks,mark=*,mark size=2.4pt,cblue,mark options={fill=white,line width=1pt},forget plot] coordinates {(87.5,6)};
\addplot[corange,line width=2.2pt,forget plot] coordinates {(80,5) (105,5)};
\addplot[only marks,mark=*,mark size=2.4pt,corange,mark options={fill=white,line width=1pt},forget plot] coordinates {(85,5)};
\addplot[corange,line width=2.2pt,forget plot] coordinates {(110,4) (136.2,4)};
\addplot[only marks,mark=*,mark size=2.4pt,corange,mark options={fill=white,line width=1pt},forget plot] coordinates {(125,4)};
\addplot[cgreen,line width=2.2pt,forget plot] coordinates {(133.8,3) (170,3)};
\addplot[only marks,mark=*,mark size=2.4pt,cgreen,mark options={fill=white,line width=1pt},forget plot] coordinates {(157.5,3)};
\node[font=\scriptsize,cmuted,anchor=west] at (axis cs:8,2) {échec (0~\% de réussite)};
\node[font=\scriptsize,cmuted,anchor=west] at (axis cs:8,1) {échec (5~\% de réussite)};
\nextgroupplot[title={S2 : radar en contrebas},ytick={8,7,6,5,4,3,2,1},yticklabels={}]
\addplot[cink,line width=2.2pt,forget plot] coordinates {(116.2,8) (191.2,8)};
\addplot[only marks,mark=*,mark size=2.4pt,cink,mark options={fill=white,line width=1pt},forget plot] coordinates {(142.5,8)};
\addplot[cblue,line width=2.2pt,forget plot] coordinates {(85,7) (110,7)};
\addplot[only marks,mark=*,mark size=2.4pt,cblue,mark options={fill=white,line width=1pt},forget plot] coordinates {(95,7)};
\addplot[cblue,line width=2.2pt,forget plot] coordinates {(88.75,6) (121.2,6)};
\addplot[only marks,mark=*,mark size=2.4pt,cblue,mark options={fill=white,line width=1pt},forget plot] coordinates {(110,6)};
\addplot[corange,line width=2.2pt,forget plot] coordinates {(108.8,5) (145,5)};
\addplot[only marks,mark=*,mark size=2.4pt,corange,mark options={fill=white,line width=1pt},forget plot] coordinates {(117.5,5)};
\addplot[corange,line width=2.2pt,forget plot] coordinates {(110,4) (140,4)};
\addplot[only marks,mark=*,mark size=2.4pt,corange,mark options={fill=white,line width=1pt},forget plot] coordinates {(127.5,4)};
\addplot[cgreen,line width=2.2pt,forget plot] coordinates {(148.8,3) (225,3)};
\addplot[only marks,mark=*,mark size=2.4pt,cgreen,mark options={fill=white,line width=1pt},forget plot] coordinates {(170,3)};
\node[font=\scriptsize,cmuted,anchor=west] at (axis cs:8,2) {échec (0~\% de réussite)};
\addplot[cmuted,line width=2.2pt,forget plot] coordinates {(142.5,1) (250,1)};
\addplot[only marks,mark=*,mark size=2.4pt,cmuted,mark options={fill=white,line width=1pt},forget plot] coordinates {(167.5,1)};
\end{groupplot}
\end{tikzpicture}

\caption{Nombre de simulations pour retrouver tout le front exact : médiane (point) et quartiles (trait) sur 20 graines, budget de 250. Deux variantes à contrainte stricte échouent.}\label{fig:l1-sims}
\end{figure}

\begin{center}
\renewcommand{\arraystretch}{1.15}
\begin{tabular}{lcccccc}
\toprule
 & \multicolumn{3}{c}{S1 (front de 4 profils)} & \multicolumn{3}{c}{S2 (front de 10 profils)} \\
 & médiane & réussite & $p$ & médiane & réussite & $p$ \\
\midrule
boucle QUBO (V2) & \textbf{75} & 100~\% & & 142,5 & 100~\% & \\
quadratique + minimum & \textbf{75} & 100~\% & 0,35 & \textbf{95} & 100~\% & $<0{,}001$ \\
krigeage + minimum & 87,5 & 100~\% & 0,18 & 110 & 100~\% & 0,002 \\
quadratique + EI & 85 & 100~\% & 0,11 & 117,5 & 100~\% & 0,03 \\
krigeage + EI & 125 & 100~\% & $<0{,}001$ & 127,5 & 100~\% & 0,06 \\
ParEGO & 157,5 & 100~\% & $<0{,}001$ & 170 & 95~\% & 0,07 \\
\midrule
quadratique + minimum, stricte & --- & 0~\% & $<0{,}001$ & --- & 0~\% & $<0{,}001$ \\
krigeage + EI, en dur & --- & 5~\% & $<0{,}001$ & 167,5 & 70~\% & 0,18 \\
\bottomrule
\end{tabular}

\smallskip
{\footnotesize Médianes en simulations distinctes sur 20 graines ; « --- » : médiane non atteinte dans le budget de 250.\\ $p$ : test de Wilcoxon apparié contre la boucle QUBO (un échec compte pour 251 simulations), sans correction pour comparaisons multiples.}
\end{center}

\begin{definition}[ : test de Wilcoxon apparié]
Les méthodes partent des mêmes points pour une graine donnée : on dispose donc, pour chaque graine, de la \emph{différence} entre les deux nombres de simulations. Le test des rangs signés de Wilcoxon demande si ces 20 différences sont réparties symétriquement autour de zéro. Il ne suppose pas qu'elles suivent une loi normale. Une petite valeur de $p$ signifie qu'un tel écart serait rare si les deux méthodes se valaient.
\end{definition}

\begin{piege}
Le tableau contient 14 tests. Au seuil de 5~\%, on attend en moyenne près d'un faux positif parmi 14 tests même si aucune méthode ne différait. Les valeurs proches du seuil (0,03 ; 0,06 ; 0,07) ne permettent donc pas de conclure. Celles inférieures à 0,002 résistent à la correction la plus sévère (Bonferroni : $0{,}05/14\approx0{,}0036$).
\end{piege}

\begin{figure}[htbp]
\centering
% figure générée par tools/make_figures.py : l1_convergence
\begin{tikzpicture}
\begin{groupplot}[group style={group size=2 by 1,horizontal sep=1.4cm},width=0.5\linewidth,height=6cm,
  xlabel={simulations distinctes},ylabel={fraction du front exact trouvée},xmin=0,xmax=250,ymin=0,ymax=1.02,
  title style={font=\small},tick label style={font=\scriptsize},label style={font=\small},
  legend style={draw=none,font=\scriptsize,at={(1.1,-0.25)},anchor=north},legend columns=3,legend cell align=left]
\nextgroupplot[title={S1 : approche de face (4 profils sur le front)}]
\addlegendimage{cink,solid,very thick}\addlegendentry{boucle QUBO (V2)}
\addlegendimage{cblue,solid,very thick}\addlegendentry{quadratique + minimum}
\addlegendimage{cblue,densely dashed,very thick}\addlegendentry{krigeage + minimum}
\addlegendimage{corange,solid,very thick}\addlegendentry{quadratique + EI}
\addlegendimage{corange,densely dashed,very thick}\addlegendentry{krigeage + EI}
\addlegendimage{cgreen,densely dashed,very thick}\addlegendentry{ParEGO}
\addplot[cink,solid,very thick,forget plot] coordinates {(0,0) (2,0) (4,0) (6,0) (8,0) (10,0) (12,0) (14,0) (16,0) (18,0) (20,0) (22,0) (24,0) (26,0) (28,0) (30,0.25) (32,0.25) (34,0.25) (36,0.25) (38,0.25) (40,0.375) (42,0.375) (44,0.375) (46,0.5) (48,0.5) (50,0.5) (52,0.5) (54,0.5) (56,0.625) (58,0.625) (60,0.75) (62,0.75) (64,0.75) (66,0.75) (68,0.75) (70,0.75) (72,0.75) (74,0.75) (76,1) (78,1) (80,1) (82,1) (84,1) (86,1) (88,1) (90,1) (92,1) (94,1) (96,1) (98,1) (100,1) (102,1) (104,1) (106,1) (108,1) (110,1) (112,1) (114,1) (116,1) (118,1) (120,1) (122,1) (124,1) (126,1) (128,1) (130,1) (132,1) (134,1) (136,1) (138,1) (140,1) (142,1) (144,1) (146,1) (148,1) (150,1) (152,1) (154,1) (156,1) (158,1) (160,1) (162,1) (164,1) (166,1) (168,1) (170,1) (172,1) (174,1) (176,1) (178,1) (180,1) (182,1) (184,1) (186,1) (188,1) (190,1) (192,1) (194,1) (196,1) (198,1) (200,1) (202,1) (204,1) (206,1) (208,1) (210,1) (212,1) (214,1) (216,1) (218,1) (220,1) (222,1) (224,1) (226,1) (228,1) (230,1) (232,1) (234,1) (236,1) (238,1) (240,1) (242,1) (244,1) (246,1) (248,1) (250,1)};
\addplot[cblue,solid,very thick,forget plot] coordinates {(0,0) (2,0) (4,0) (6,0) (8,0) (10,0) (12,0) (14,0) (16,0) (18,0) (20,0) (22,0) (24,0) (26,0) (28,0) (30,0) (32,0) (34,0) (36,0.125) (38,0.125) (40,0.25) (42,0.25) (44,0.25) (46,0.5) (48,0.5) (50,0.5) (52,0.5) (54,0.5) (56,0.5) (58,0.5) (60,0.5) (62,0.5) (64,0.5) (66,0.75) (68,0.75) (70,0.75) (72,0.75) (74,0.75) (76,1) (78,1) (80,1) (82,1) (84,1) (86,1) (88,1) (90,1) (92,1) (94,1) (96,1) (98,1) (100,1) (102,1) (104,1) (106,1) (108,1) (110,1) (112,1) (114,1) (116,1) (118,1) (120,1) (122,1) (124,1) (126,1) (128,1) (130,1) (132,1) (134,1) (136,1) (138,1) (140,1) (142,1) (144,1) (146,1) (148,1) (150,1) (152,1) (154,1) (156,1) (158,1) (160,1) (162,1) (164,1) (166,1) (168,1) (170,1) (172,1) (174,1) (176,1) (178,1) (180,1) (182,1) (184,1) (186,1) (188,1) (190,1) (192,1) (194,1) (196,1) (198,1) (200,1) (202,1) (204,1) (206,1) (208,1) (210,1) (212,1) (214,1) (216,1) (218,1) (220,1) (222,1) (224,1) (226,1) (228,1) (230,1) (232,1) (234,1) (236,1) (238,1) (240,1) (242,1) (244,1) (246,1) (248,1) (250,1)};
\addplot[cblue,densely dashed,very thick,forget plot] coordinates {(0,0) (2,0) (4,0) (6,0) (8,0) (10,0) (12,0) (14,0) (16,0) (18,0) (20,0) (22,0) (24,0) (26,0) (28,0) (30,0.25) (32,0.25) (34,0.25) (36,0.25) (38,0.25) (40,0.5) (42,0.5) (44,0.5) (46,0.5) (48,0.5) (50,0.5) (52,0.5) (54,0.5) (56,0.5) (58,0.5) (60,0.5) (62,0.5) (64,0.5) (66,0.5) (68,0.5) (70,0.75) (72,0.75) (74,0.75) (76,0.75) (78,0.75) (80,0.75) (82,0.75) (84,0.75) (86,0.875) (88,0.875) (90,1) (92,1) (94,1) (96,1) (98,1) (100,1) (102,1) (104,1) (106,1) (108,1) (110,1) (112,1) (114,1) (116,1) (118,1) (120,1) (122,1) (124,1) (126,1) (128,1) (130,1) (132,1) (134,1) (136,1) (138,1) (140,1) (142,1) (144,1) (146,1) (148,1) (150,1) (152,1) (154,1) (156,1) (158,1) (160,1) (162,1) (164,1) (166,1) (168,1) (170,1) (172,1) (174,1) (176,1) (178,1) (180,1) (182,1) (184,1) (186,1) (188,1) (190,1) (192,1) (194,1) (196,1) (198,1) (200,1) (202,1) (204,1) (206,1) (208,1) (210,1) (212,1) (214,1) (216,1) (218,1) (220,1) (222,1) (224,1) (226,1) (228,1) (230,1) (232,1) (234,1) (236,1) (238,1) (240,1) (242,1) (244,1) (246,1) (248,1) (250,1)};
\addplot[corange,solid,very thick,forget plot] coordinates {(0,0) (2,0) (4,0) (6,0) (8,0) (10,0) (12,0) (14,0) (16,0) (18,0) (20,0) (22,0) (24,0) (26,0) (28,0) (30,0) (32,0) (34,0) (36,0.25) (38,0.25) (40,0.25) (42,0.25) (44,0.25) (46,0.5) (48,0.5) (50,0.5) (52,0.5) (54,0.5) (56,0.5) (58,0.5) (60,0.5) (62,0.5) (64,0.5) (66,0.625) (68,0.625) (70,0.75) (72,0.75) (74,0.75) (76,0.75) (78,0.75) (80,0.75) (82,0.75) (84,0.75) (86,1) (88,1) (90,1) (92,1) (94,1) (96,1) (98,1) (100,1) (102,1) (104,1) (106,1) (108,1) (110,1) (112,1) (114,1) (116,1) (118,1) (120,1) (122,1) (124,1) (126,1) (128,1) (130,1) (132,1) (134,1) (136,1) (138,1) (140,1) (142,1) (144,1) (146,1) (148,1) (150,1) (152,1) (154,1) (156,1) (158,1) (160,1) (162,1) (164,1) (166,1) (168,1) (170,1) (172,1) (174,1) (176,1) (178,1) (180,1) (182,1) (184,1) (186,1) (188,1) (190,1) (192,1) (194,1) (196,1) (198,1) (200,1) (202,1) (204,1) (206,1) (208,1) (210,1) (212,1) (214,1) (216,1) (218,1) (220,1) (222,1) (224,1) (226,1) (228,1) (230,1) (232,1) (234,1) (236,1) (238,1) (240,1) (242,1) (244,1) (246,1) (248,1) (250,1)};
\addplot[corange,densely dashed,very thick,forget plot] coordinates {(0,0) (2,0) (4,0) (6,0) (8,0) (10,0) (12,0) (14,0) (16,0) (18,0) (20,0) (22,0) (24,0) (26,0) (28,0) (30,0) (32,0) (34,0) (36,0.125) (38,0.125) (40,0.125) (42,0.125) (44,0.125) (46,0.25) (48,0.25) (50,0.25) (52,0.25) (54,0.25) (56,0.5) (58,0.5) (60,0.5) (62,0.5) (64,0.5) (66,0.5) (68,0.5) (70,0.5) (72,0.5) (74,0.5) (76,0.5) (78,0.5) (80,0.5) (82,0.5) (84,0.5) (86,0.5) (88,0.5) (90,0.5) (92,0.5) (94,0.5) (96,0.75) (98,0.75) (100,0.75) (102,0.75) (104,0.75) (106,0.75) (108,0.75) (110,0.75) (112,0.75) (114,0.75) (116,0.75) (118,0.75) (120,0.75) (122,0.75) (124,0.75) (126,1) (128,1) (130,1) (132,1) (134,1) (136,1) (138,1) (140,1) (142,1) (144,1) (146,1) (148,1) (150,1) (152,1) (154,1) (156,1) (158,1) (160,1) (162,1) (164,1) (166,1) (168,1) (170,1) (172,1) (174,1) (176,1) (178,1) (180,1) (182,1) (184,1) (186,1) (188,1) (190,1) (192,1) (194,1) (196,1) (198,1) (200,1) (202,1) (204,1) (206,1) (208,1) (210,1) (212,1) (214,1) (216,1) (218,1) (220,1) (222,1) (224,1) (226,1) (228,1) (230,1) (232,1) (234,1) (236,1) (238,1) (240,1) (242,1) (244,1) (246,1) (248,1) (250,1)};
\addplot[cgreen,densely dashed,very thick,forget plot] coordinates {(0,0) (2,0) (4,0) (6,0) (8,0) (10,0) (12,0) (14,0) (16,0) (18,0) (20,0) (22,0) (24,0) (26,0) (28,0) (30,0.125) (32,0.125) (34,0.125) (36,0.25) (38,0.25) (40,0.25) (42,0.25) (44,0.25) (46,0.25) (48,0.25) (50,0.5) (52,0.5) (54,0.5) (56,0.5) (58,0.5) (60,0.5) (62,0.5) (64,0.5) (66,0.5) (68,0.5) (70,0.5) (72,0.5) (74,0.5) (76,0.5) (78,0.5) (80,0.5) (82,0.5) (84,0.5) (86,0.5) (88,0.5) (90,0.5) (92,0.5) (94,0.5) (96,0.5) (98,0.5) (100,0.5) (102,0.5) (104,0.5) (106,0.5) (108,0.5) (110,0.5) (112,0.5) (114,0.5) (116,0.5) (118,0.5) (120,0.75) (122,0.75) (124,0.75) (126,0.75) (128,0.75) (130,0.75) (132,0.75) (134,0.75) (136,0.75) (138,0.75) (140,0.75) (142,0.75) (144,0.75) (146,0.75) (148,0.75) (150,0.75) (152,0.75) (154,0.75) (156,0.875) (158,0.875) (160,1) (162,1) (164,1) (166,1) (168,1) (170,1) (172,1) (174,1) (176,1) (178,1) (180,1) (182,1) (184,1) (186,1) (188,1) (190,1) (192,1) (194,1) (196,1) (198,1) (200,1) (202,1) (204,1) (206,1) (208,1) (210,1) (212,1) (214,1) (216,1) (218,1) (220,1) (222,1) (224,1) (226,1) (228,1) (230,1) (232,1) (234,1) (236,1) (238,1) (240,1) (242,1) (244,1) (246,1) (248,1) (250,1)};
\nextgroupplot[title={S2 : radar en contrebas (10 profils sur le front)}]
\addplot[cink,solid,very thick,forget plot] coordinates {(0,0) (2,0) (4,0) (6,0) (8,0) (10,0) (12,0) (14,0) (16,0) (18,0) (20,0) (22,0) (24,0) (26,0) (28,0) (30,0.1) (32,0.1) (34,0.1) (36,0.1) (38,0.1) (40,0.1) (42,0.1) (44,0.1) (46,0.2) (48,0.2) (50,0.3) (52,0.3) (54,0.3) (56,0.3) (58,0.3) (60,0.4) (62,0.4) (64,0.4) (66,0.5) (68,0.5) (70,0.55) (72,0.55) (74,0.55) (76,0.6) (78,0.6) (80,0.7) (82,0.7) (84,0.7) (86,0.7) (88,0.7) (90,0.75) (92,0.75) (94,0.75) (96,0.8) (98,0.8) (100,0.85) (102,0.85) (104,0.85) (106,0.85) (108,0.85) (110,0.85) (112,0.85) (114,0.85) (116,0.9) (118,0.9) (120,0.9) (122,0.9) (124,0.9) (126,0.9) (128,0.9) (130,0.9) (132,0.9) (134,0.9) (136,0.9) (138,0.9) (140,0.95) (142,0.95) (144,0.95) (146,1) (148,1) (150,1) (152,1) (154,1) (156,1) (158,1) (160,1) (162,1) (164,1) (166,1) (168,1) (170,1) (172,1) (174,1) (176,1) (178,1) (180,1) (182,1) (184,1) (186,1) (188,1) (190,1) (192,1) (194,1) (196,1) (198,1) (200,1) (202,1) (204,1) (206,1) (208,1) (210,1) (212,1) (214,1) (216,1) (218,1) (220,1) (222,1) (224,1) (226,1) (228,1) (230,1) (232,1) (234,1) (236,1) (238,1) (240,1) (242,1) (244,1) (246,1) (248,1) (250,1)};
\addplot[cblue,solid,very thick,forget plot] coordinates {(0,0) (2,0) (4,0) (6,0) (8,0) (10,0) (12,0) (14,0) (16,0) (18,0) (20,0) (22,0) (24,0) (26,0.05) (28,0.05) (30,0.2) (32,0.2) (34,0.2) (36,0.2) (38,0.2) (40,0.2) (42,0.2) (44,0.2) (46,0.3) (48,0.3) (50,0.45) (52,0.45) (54,0.45) (56,0.5) (58,0.5) (60,0.6) (62,0.6) (64,0.6) (66,0.7) (68,0.7) (70,0.7) (72,0.7) (74,0.7) (76,0.8) (78,0.8) (80,0.9) (82,0.9) (84,0.9) (86,0.9) (88,0.9) (90,0.9) (92,0.9) (94,0.9) (96,1) (98,1) (100,1) (102,1) (104,1) (106,1) (108,1) (110,1) (112,1) (114,1) (116,1) (118,1) (120,1) (122,1) (124,1) (126,1) (128,1) (130,1) (132,1) (134,1) (136,1) (138,1) (140,1) (142,1) (144,1) (146,1) (148,1) (150,1) (152,1) (154,1) (156,1) (158,1) (160,1) (162,1) (164,1) (166,1) (168,1) (170,1) (172,1) (174,1) (176,1) (178,1) (180,1) (182,1) (184,1) (186,1) (188,1) (190,1) (192,1) (194,1) (196,1) (198,1) (200,1) (202,1) (204,1) (206,1) (208,1) (210,1) (212,1) (214,1) (216,1) (218,1) (220,1) (222,1) (224,1) (226,1) (228,1) (230,1) (232,1) (234,1) (236,1) (238,1) (240,1) (242,1) (244,1) (246,1) (248,1) (250,1)};
\addplot[cblue,densely dashed,very thick,forget plot] coordinates {(0,0) (2,0) (4,0) (6,0) (8,0) (10,0) (12,0) (14,0) (16,0) (18,0) (20,0) (22,0) (24,0) (26,0.1) (28,0.1) (30,0.1) (32,0.1) (34,0.1) (36,0.2) (38,0.2) (40,0.25) (42,0.25) (44,0.25) (46,0.3) (48,0.3) (50,0.4) (52,0.4) (54,0.4) (56,0.4) (58,0.4) (60,0.5) (62,0.5) (64,0.5) (66,0.6) (68,0.6) (70,0.65) (72,0.65) (74,0.65) (76,0.7) (78,0.7) (80,0.8) (82,0.8) (84,0.8) (86,0.8) (88,0.8) (90,0.9) (92,0.9) (94,0.9) (96,0.9) (98,0.9) (100,0.9) (102,0.9) (104,0.9) (106,0.9) (108,0.9) (110,1) (112,1) (114,1) (116,1) (118,1) (120,1) (122,1) (124,1) (126,1) (128,1) (130,1) (132,1) (134,1) (136,1) (138,1) (140,1) (142,1) (144,1) (146,1) (148,1) (150,1) (152,1) (154,1) (156,1) (158,1) (160,1) (162,1) (164,1) (166,1) (168,1) (170,1) (172,1) (174,1) (176,1) (178,1) (180,1) (182,1) (184,1) (186,1) (188,1) (190,1) (192,1) (194,1) (196,1) (198,1) (200,1) (202,1) (204,1) (206,1) (208,1) (210,1) (212,1) (214,1) (216,1) (218,1) (220,1) (222,1) (224,1) (226,1) (228,1) (230,1) (232,1) (234,1) (236,1) (238,1) (240,1) (242,1) (244,1) (246,1) (248,1) (250,1)};
\addplot[corange,solid,very thick,forget plot] coordinates {(0,0) (2,0) (4,0) (6,0) (8,0) (10,0) (12,0) (14,0) (16,0) (18,0) (20,0) (22,0) (24,0) (26,0.1) (28,0.1) (30,0.1) (32,0.1) (34,0.1) (36,0.2) (38,0.2) (40,0.2) (42,0.2) (44,0.2) (46,0.3) (48,0.3) (50,0.35) (52,0.35) (54,0.35) (56,0.4) (58,0.4) (60,0.5) (62,0.5) (64,0.5) (66,0.5) (68,0.5) (70,0.5) (72,0.5) (74,0.5) (76,0.7) (78,0.7) (80,0.7) (82,0.7) (84,0.7) (86,0.7) (88,0.7) (90,0.8) (92,0.8) (94,0.8) (96,0.8) (98,0.8) (100,0.9) (102,0.9) (104,0.9) (106,0.9) (108,0.9) (110,0.9) (112,0.9) (114,0.9) (116,0.95) (118,0.95) (120,1) (122,1) (124,1) (126,1) (128,1) (130,1) (132,1) (134,1) (136,1) (138,1) (140,1) (142,1) (144,1) (146,1) (148,1) (150,1) (152,1) (154,1) (156,1) (158,1) (160,1) (162,1) (164,1) (166,1) (168,1) (170,1) (172,1) (174,1) (176,1) (178,1) (180,1) (182,1) (184,1) (186,1) (188,1) (190,1) (192,1) (194,1) (196,1) (198,1) (200,1) (202,1) (204,1) (206,1) (208,1) (210,1) (212,1) (214,1) (216,1) (218,1) (220,1) (222,1) (224,1) (226,1) (228,1) (230,1) (232,1) (234,1) (236,1) (238,1) (240,1) (242,1) (244,1) (246,1) (248,1) (250,1)};
\addplot[corange,densely dashed,very thick,forget plot] coordinates {(0,0) (2,0) (4,0) (6,0) (8,0) (10,0) (12,0) (14,0) (16,0) (18,0) (20,0) (22,0) (24,0) (26,0.1) (28,0.1) (30,0.1) (32,0.1) (34,0.1) (36,0.15) (38,0.15) (40,0.2) (42,0.2) (44,0.2) (46,0.3) (48,0.3) (50,0.3) (52,0.3) (54,0.3) (56,0.4) (58,0.4) (60,0.5) (62,0.5) (64,0.5) (66,0.6) (68,0.6) (70,0.6) (72,0.6) (74,0.6) (76,0.7) (78,0.7) (80,0.7) (82,0.7) (84,0.7) (86,0.7) (88,0.7) (90,0.8) (92,0.8) (94,0.8) (96,0.8) (98,0.8) (100,0.8) (102,0.8) (104,0.8) (106,0.85) (108,0.85) (110,0.9) (112,0.9) (114,0.9) (116,0.9) (118,0.9) (120,0.9) (122,0.9) (124,0.9) (126,0.95) (128,0.95) (130,1) (132,1) (134,1) (136,1) (138,1) (140,1) (142,1) (144,1) (146,1) (148,1) (150,1) (152,1) (154,1) (156,1) (158,1) (160,1) (162,1) (164,1) (166,1) (168,1) (170,1) (172,1) (174,1) (176,1) (178,1) (180,1) (182,1) (184,1) (186,1) (188,1) (190,1) (192,1) (194,1) (196,1) (198,1) (200,1) (202,1) (204,1) (206,1) (208,1) (210,1) (212,1) (214,1) (216,1) (218,1) (220,1) (222,1) (224,1) (226,1) (228,1) (230,1) (232,1) (234,1) (236,1) (238,1) (240,1) (242,1) (244,1) (246,1) (248,1) (250,1)};
\addplot[cgreen,densely dashed,very thick,forget plot] coordinates {(0,0) (2,0) (4,0) (6,0) (8,0) (10,0) (12,0) (14,0) (16,0) (18,0) (20,0) (22,0) (24,0) (26,0) (28,0) (30,0.1) (32,0.1) (34,0.1) (36,0.15) (38,0.15) (40,0.2) (42,0.2) (44,0.2) (46,0.2) (48,0.2) (50,0.25) (52,0.25) (54,0.25) (56,0.3) (58,0.3) (60,0.35) (62,0.35) (64,0.35) (66,0.4) (68,0.4) (70,0.4) (72,0.4) (74,0.4) (76,0.45) (78,0.45) (80,0.5) (82,0.5) (84,0.5) (86,0.5) (88,0.5) (90,0.5) (92,0.5) (94,0.5) (96,0.6) (98,0.6) (100,0.6) (102,0.6) (104,0.6) (106,0.65) (108,0.65) (110,0.7) (112,0.7) (114,0.7) (116,0.7) (118,0.7) (120,0.75) (122,0.75) (124,0.75) (126,0.8) (128,0.8) (130,0.8) (132,0.8) (134,0.8) (136,0.85) (138,0.85) (140,0.9) (142,0.9) (144,0.9) (146,0.9) (148,0.9) (150,0.9) (152,0.9) (154,0.9) (156,0.9) (158,0.9) (160,0.9) (162,0.9) (164,0.9) (166,0.9) (168,0.9) (170,1) (172,1) (174,1) (176,1) (178,1) (180,1) (182,1) (184,1) (186,1) (188,1) (190,1) (192,1) (194,1) (196,1) (198,1) (200,1) (202,1) (204,1) (206,1) (208,1) (210,1) (212,1) (214,1) (216,1) (218,1) (220,1) (222,1) (224,1) (226,1) (228,1) (230,1) (232,1) (234,1) (236,1) (238,1) (240,1) (242,1) (244,1) (246,1) (248,1) (250,1)};
\end{groupplot}
\end{tikzpicture}

\caption{Fraction du front exact trouvée selon le nombre de simulations (médiane sur 20 graines).}\label{fig:l1-conv}
\end{figure}

\section{Avec un arrêt réaliste}
Sans connaître le front exact, on s'arrête quand l'hypervolume ne progresse plus depuis 3 itérations (même règle que la V2). Médianes au moment de l'arrêt :
\begin{center}
\begin{tabular}{lcccccc}
\toprule
 & \multicolumn{3}{c}{S1} & \multicolumn{3}{c}{S2} \\
 & simulations & front & HV & simulations & front & HV \\
\midrule
boucle QUBO (V2) & 52,5 & 50~\% & 0,943 & 67,5 & 55~\% & 0,962 \\
quadratique + minimum & 65 & 50~\% & 0,937 & 100 & 100~\% & 1,000 \\
krigeage + minimum & 50 & 50~\% & 0,937 & 95 & 90~\% & 0,994 \\
quadratique + EI & 57,5 & 50~\% & 0,937 & 70 & 55~\% & 0,968 \\
krigeage + EI & 55 & 50~\% & 0,937 & 100 & 90~\% & 0,987 \\
ParEGO & 57,5 & 50~\% & 0,937 & 100 & 55~\% & 0,985 \\
\bottomrule
\end{tabular}
\end{center}
Dans S1, toutes les méthodes s'arrêtent au même point : la moitié du front et 94~\% de l'hypervolume, en 50 à 65 simulations. Dans S2, la boucle QUBO s'arrête plus tôt (67,5 simulations) mais avec 55~\% du front ; trois variantes continuent jusqu'à environ 100 simulations et s'arrêtent avec 90 à 100~\% du front.

\section{Le coût de calcul}
\begin{center}
\begin{tabular}{lcc}
\toprule
 & S1 & S2 \\
\midrule
boucle QUBO (V2), énumération de $2^{14}$ états & 2,8~s & 6,0~s \\
quadratique + minimum & 0,9~s & 0,6~s \\
quadratique + EI & 0,4~s & 0,8~s \\
krigeage + minimum & 24~s & 27~s \\
krigeage + EI & 34~s & 32~s \\
ParEGO & 84~s & 104~s \\
\bottomrule
\end{tabular}

\smallskip
{\footnotesize Temps moyen par graine, hors simulations (qui sont lues dans la table), sur un ordinateur portable.}
\end{center}
Ces temps restent négligeables devant une simulation réelle (XFOIL : environ 3~s par profil ; un calcul de mécanique des fluides 3D : des heures). Ils ne deviennent un critère que si deux méthodes demandent le même nombre de simulations.

\section{Lecture}
\paragraph{1. H1 n'est pas confirmée.} L'optimisation bayésienne (krigeage + EI) retrouve le front à chaque fois, mais elle ne bat pas la boucle QUBO (figures~\ref{fig:l1-sims} et~\ref{fig:l1-conv}). Dans S1, elle demande nettement plus de simulations (125 contre 75, $p<0{,}001$). Dans S2, l'écart en sa faveur (127,5 contre 142,5) n'est pas significatif.

\paragraph{2. Le krigeage n'apporte rien sur ce problème.} À sélection égale, le modèle quadratique fait aussi bien ou mieux que le krigeage dans les quatre comparaisons (75 contre 87,5 et 85 contre 125 dans S1 ; 95 contre 110 et 117,5 contre 127,5 dans S2), pour un coût de calcul 25 à 30 fois plus faible. Le chapitre~\ref{ch:diagnostic} regarde pourquoi.

\paragraph{3. L'amélioration espérée non plus.} À modèle égal, le minimum de la prédiction fait aussi bien ou mieux que l'amélioration espérée. L'exploration que celle-ci ajoute n'est pas utile ici, parce que la pénalité lâche sur la contrainte en fournit déjà (chapitre~\ref{ch:diagnostic}). Réserve : les deux sélections ne traitent pas la contrainte de la même façon (pénalité de la boucle QUBO pour le minimum, probabilité de faisabilité pour l'amélioration espérée). L'écart mesuré mêle donc les deux effets.

\paragraph{4. Le traitement de la contrainte domine.} Les deux variantes à contrainte stricte échouent, quel que soit le modèle : 0~\% et 5~\% de réussite dans S1. Aucun autre choix du tableau n'a un effet de cette taille.

\paragraph{5. Dans S2, le contrôle bat la boucle QUBO.} La variante « quadratique + minimum » reprend le modèle, la sélection et la pénalité de la boucle QUBO, sans encodage binaire. Elle devait se comporter comme elle. C'est le cas dans S1 (75 et 75), mais dans S2 elle retrouve le front en 95 simulations contre 142,5 ($p<0{,}001$), et s'arrête avec tout le front au lieu de 55~\%.

\begin{piege}
Ce dernier résultat ne permet pas encore de dire que « l'encodage en QUBO coûte ». Le contrôle diffère de la boucle QUBO par trois choses à la fois : l'absence de qubits d'écart (la variable d'écart n'a que 8 valeurs, ce qui biaise la pénalité), la régularisation du modèle quadratique (bayésienne ici, ridge fixe à $10^{-6}$ dans la boucle), et l'épaisseur minimale traitée par filtre au lieu d'une pénalité. On ne sait pas laquelle explique l'écart. C'est la première piste de ce chapitre.
\end{piege}

\begin{resultat}
\begin{itemize}
\item L'optimisation bayésienne n'égale la boucle QUBO que dans un scénario sur deux : H1 n'est pas confirmée.
\item Ni le krigeage ni l'amélioration espérée n'améliorent la boucle. Un modèle quadratique et le minimum de sa prédiction suffisent.
\item Le choix qui compte est le traitement de la contrainte de manœuvre.
\end{itemize}
\end{resultat}

\exercices
\begin{exercice}
Dans S1, le krigeage + EI demande 125 simulations contre 75 pour la boucle QUBO, mais avec l'arrêt réaliste les deux s'arrêtent vers 55 simulations avec la moitié du front. Ces deux constats sont-ils contradictoires ?
\end{exercice}
\begin{exercice}
Sur 20 graines, une méthode A retrouve le front plus vite que B pour 15 graines, plus lentement pour 5. Sans calcul de rangs, quelle est la probabilité d'un résultat au moins aussi déséquilibré si A et B se valaient (test du signe) ? Conclure au seuil de 5~\%.
\end{exercice}

\begin{pistes}
\begin{itemize}
\item \textbf{Identifier la cause de l'écart dans S2} entre le contrôle et la boucle QUBO, en changeant un seul des trois éléments à la fois (qubits d'écart, régularisation, filtre d'épaisseur).
\item Remplacer le critère « tout le front exact », fragile pour des fronts de 4 et 10 profils, par l'hypervolume ou la distance au front en fonction du budget.
\item Corriger les valeurs de $p$ pour comparaisons multiples (Holm), et ajouter des intervalles de confiance par rééchantillonnage sur les médianes.
\end{itemize}
\end{pistes}

\chapter{Ce qui explique les résultats : la contrainte de manœuvre}\label{ch:diagnostic}
\iterbanner{V3}

Le chapitre précédent laisse deux questions. Pourquoi les variantes à contrainte stricte échouent-elles ? Et pourquoi le krigeage, plus précis sur un cas jouet, n'aide-t-il pas ? Les deux diagnostics de \code{optim/diagnostic\_l1.py} y répondent.

\section{Le front est au bord de la contrainte}
Les 4 profils du front de S1 ont un $C_{L,\max}$ de 1,20 à 1,34, pour une exigence de 1,2. Ce n'est pas un hasard : un nez aigu réduit la SER et décroche tôt, donc les meilleurs compromis sont ceux qui tiennent la manœuvre de justesse (cours d'optimisation, chapitre V2).

Or les modèles de $C_{L,\max}$ sont imprécis à cet endroit. Entraînons les trois modèles sur \emph{toute} la table (le cas le plus favorable) et regardons ce qu'ils prédisent pour ces 4 profils :
\begin{center}
\renewcommand{\arraystretch}{1.15}
\begin{tabular}{ccccc}
\toprule
profil $(m;\ p;\ t;\ I)$ & simulé & modèle affine & modèle quadratique & krigeage \\
\midrule
2~\% ; 0,45 ; 8~\% ; 4 & 1,30 & 1,24 & $1{,}11\pm0{,}15$ & $1{,}17\pm0{,}02$ \\
2~\% ; 0,55 ; 8~\% ; 4 & 1,32 & 1,23 & $1{,}13\pm0{,}15$ & $1{,}16\pm0{,}03$ \\
4~\% ; 0,55 ; 12~\% ; 2 & 1,20 & 1,26 & $1{,}12\pm0{,}15$ & $1{,}04\pm0{,}03$ \\
5~\% ; 0,45 ; 10~\% ; 2 & 1,34 & 1,33 & $1{,}06\pm0{,}15$ & $1{,}06\pm0{,}03$ \\
\midrule
verdict sur le front & faisables & faisables & infaisables & infaisables \\
pénalité subie & & 0 & 0,69 à 0,83 & 0,92 à 1,00 \\
erreur moyenne sur la grille & & 0,19 & 0,14 & \\
\bottomrule
\end{tabular}

\smallskip
{\footnotesize Pénalité : $75\max(1{,}2-\widehat{C_{L,\max}},0)^2$ pour le modèle affine (boucle QUBO), $1-\mathrm{PoF}$ pour les deux autres ; à comparer à l'étendue des objectifs normalisés, de l'ordre de 1.}
\end{center}

\begin{resultat}
Le modèle affine est le plus grossier des trois (erreur moyenne de 0,19 contre 0,14 pour le quadratique), mais c'est le seul qui laisse passer les quatre profils du front. Les deux modèles plus précis en moyenne les prédisent infaisables ; le krigeage en est même presque certain. Une méthode qui leur fait confiance ne simule jamais ces profils, donc ne les trouve jamais.
\end{resultat}

\begin{projet}
La boucle QUBO utilise un modèle affine de $C_{L,\max}$ par \emph{nécessité} : une contrainte doit être de degré~1 pour que sa pénalité reste quadratique (cours d'optimisation, chapitre sur la boucle). Cette limitation de la formulation se trouve être un avantage ici : le modèle affine est optimiste près du front (il prédit 84~\% de profils faisables sur la grille, contre 77~\% en réalité). On parlera dans la suite de traitement \emph{lâche} de la contrainte.
\end{projet}

\section{Premier diagnostic : combien de profils infaisables faut-il oser simuler ?}
Sur la grille, 23~\% des profils violent la contrainte de manœuvre. Pour chaque méthode, on mesure la part des profils \emph{proposés} (hors les 10 initiaux) qui la violent, et le nombre de simulations pour retrouver le front. Pour la boucle QUBO, on fait varier le poids de la pénalité (25 dans la V2).

\begin{figure}[htbp]
\centering
% figure générée par tools/make_figures.py : diag_contrainte
\begin{tikzpicture}
\begin{groupplot}[group style={group size=2 by 1,horizontal sep=1.1cm},width=0.5\linewidth,height=6.2cm,
  xlabel={profils proposés infaisables (\%)},
  xmin=0,xmax=85,ymin=0,ymax=290,ytick={0,50,100,150,200,250,270},yticklabels={0,50,100,150,200,250,échec},
  title style={font=\small},tick label style={font=\scriptsize},label style={font=\small}]
\nextgroupplot[title={S1 : approche de face},ylabel={simulations pour le front complet}]
\draw[cmuted,dashed] (axis cs:22.8,0) -- (axis cs:22.8,290);
\addplot[only marks,mark=square*,mark size=2.6pt,cink,forget plot] coordinates {(54.8,85)};
\node[font=\scriptsize,anchor=south west,inner sep=2pt] at (axis cs:54.8,85) {1};
\addplot[only marks,mark=square*,mark size=2.6pt,cink,forget plot] coordinates {(45.4,75)};
\node[font=\scriptsize,anchor=south west,inner sep=2pt] at (axis cs:45.4,75) {2};
\addplot[only marks,mark=square*,mark size=2.6pt,cink,forget plot] coordinates {(32.3,107.5)};
\node[font=\scriptsize,anchor=south west,inner sep=2pt] at (axis cs:32.3,107.5) {3};
\addplot[only marks,mark=square*,mark size=2.6pt,cink,forget plot] coordinates {(70.8,112.5)};
\node[font=\scriptsize,anchor=south west,inner sep=2pt] at (axis cs:70.8,112.5) {4};
\addplot[only marks,mark=*,mark size=2.6pt,cblue,forget plot] coordinates {(41.3,75)};
\node[font=\scriptsize,anchor=south west,inner sep=2pt] at (axis cs:41.3,75) {5};
\addplot[only marks,mark=*,mark size=2.6pt,cblue,forget plot] coordinates {(8.12,270)};
\node[font=\scriptsize,anchor=south west,inner sep=2pt] at (axis cs:8.12,270) {6};
\addplot[only marks,mark=*,mark size=2.6pt,cblue,forget plot] coordinates {(70.1,120)};
\node[font=\scriptsize,anchor=south west,inner sep=2pt] at (axis cs:70.1,120) {7};
\addplot[only marks,mark=triangle*,mark size=2.6pt,corange,forget plot] coordinates {(58.3,85)};
\node[font=\scriptsize,anchor=south west,inner sep=2pt] at (axis cs:58.3,85) {8};
\addplot[only marks,mark=triangle*,mark size=2.6pt,corange,forget plot] coordinates {(77.9,115)};
\node[font=\scriptsize,anchor=south west,inner sep=2pt] at (axis cs:77.9,115) {9};
\nextgroupplot[title={S2 : radar en contrebas}]
\draw[cmuted,dashed] (axis cs:22.8,0) -- (axis cs:22.8,290);
\addplot[only marks,mark=square*,mark size=2.6pt,cink,forget plot] coordinates {(43.5,110)};
\node[font=\scriptsize,anchor=south west,inner sep=2pt] at (axis cs:43.5,110) {1};
\addplot[only marks,mark=square*,mark size=2.6pt,cink,forget plot] coordinates {(33.2,142.5)};
\node[font=\scriptsize,anchor=south west,inner sep=2pt] at (axis cs:33.2,142.5) {2};
\addplot[only marks,mark=square*,mark size=2.6pt,cink,forget plot] coordinates {(23.6,270)};
\node[font=\scriptsize,anchor=south west,inner sep=2pt] at (axis cs:23.6,270) {3};
\addplot[only marks,mark=square*,mark size=2.6pt,cink,forget plot] coordinates {(52.1,147.5)};
\node[font=\scriptsize,anchor=south west,inner sep=2pt] at (axis cs:52.1,147.5) {4};
\addplot[only marks,mark=*,mark size=2.6pt,cblue,forget plot] coordinates {(38.4,95)};
\node[font=\scriptsize,anchor=south west,inner sep=2pt] at (axis cs:38.4,95) {5};
\addplot[only marks,mark=*,mark size=2.6pt,cblue,forget plot] coordinates {(8.25,270)};
\node[font=\scriptsize,anchor=south west,inner sep=2pt] at (axis cs:8.25,270) {6};
\addplot[only marks,mark=*,mark size=2.6pt,cblue,forget plot] coordinates {(56,125)};
\node[font=\scriptsize,anchor=south west,inner sep=2pt] at (axis cs:56,125) {7};
\addplot[only marks,mark=triangle*,mark size=2.6pt,corange,forget plot] coordinates {(40.7,117.5)};
\node[font=\scriptsize,anchor=south west,inner sep=2pt] at (axis cs:40.7,117.5) {8};
\addplot[only marks,mark=triangle*,mark size=2.6pt,corange,forget plot] coordinates {(53.3,152.5)};
\node[font=\scriptsize,anchor=south west,inner sep=2pt] at (axis cs:53.3,152.5) {9};
\end{groupplot}
\end{tikzpicture}

\caption{Part des profils proposés qui violent la contrainte de manœuvre, et nombre médian de simulations pour retrouver le front (20 graines). Carrés : boucle QUBO, pénalité de poids 2,5 (1), 25 (2, réglage V2), 250 (3), sans contrainte (4). Disques : quadratique + minimum, avec la pénalité de la boucle QUBO (5), la pénalité probabiliste stricte (6), sans contrainte (7). Triangles : quadratique + EI, avec probabilité de faisabilité (8), sans contrainte (9). Tirets : part des profils infaisables sur la grille (23~\%).}\label{fig:diag-contrainte}
\end{figure}

\begin{center}
\small
\renewcommand{\arraystretch}{1.15}
\setlength{\tabcolsep}{4pt}
\begin{tabular}{lcccccc}
\toprule
 & \multicolumn{3}{c}{S1} & \multicolumn{3}{c}{S2} \\
 & violent & médiane & réussite & violent & médiane & réussite \\
\midrule
boucle QUBO, sans contrainte & 71~\% & 112,5 & 100~\% & 52~\% & 147,5 & 100~\% \\
boucle QUBO, poids 2,5 & 55~\% & 85 & 100~\% & 44~\% & \textbf{110} & 100~\% \\
boucle QUBO, poids 25 (V2) & 45~\% & \textbf{75} & 100~\% & 33~\% & 142,5 & 100~\% \\
boucle QUBO, poids 250 & 32~\% & 107,5 & 75~\% & 24~\% & --- & 25~\% \\
\midrule
quadratique + minimum, sans contrainte & 70~\% & 120 & 100~\% & 56~\% & 125 & 100~\% \\
quadratique + minimum, pénalité QUBO & 41~\% & \textbf{75} & 100~\% & 38~\% & \textbf{95} & 100~\% \\
quadratique + minimum, stricte & 8~\% & --- & 0~\% & 8~\% & --- & 0~\% \\
\midrule
quadratique + EI, sans contrainte & 78~\% & 115 & 100~\% & 53~\% & 152,5 & 100~\% \\
quadratique + EI, probabilité de faisabilité & 58~\% & 85 & 100~\% & 41~\% & 117,5 & 100~\% \\
\bottomrule
\end{tabular}
\end{center}

\paragraph{Lecture.}
\begin{itemize}
\item \textbf{Trop strict, on échoue.} La variante qui ne propose que 8~\% de profils infaisables ne retrouve jamais le front. La boucle QUBO avec une pénalité dix fois plus forte tombe à 75~\% de réussite dans S1 et 25~\% dans S2.
\item \textbf{Sans contrainte, on réussit, mais plus lentement.} Ignorer la manœuvre dans le modèle coûte 30 à 60~\% de simulations en plus dans cinq cas sur six (la boucle QUBO dans S2 fait exception : 147,5 contre 142,5) : on gaspille des simulations sur des profils nettement infaisables.
\item \textbf{Les méthodes qui réussissent le mieux proposent 33 à 58~\% de profils infaisables}, soit une fois et demie à deux fois et demie leur part sur la grille. Elles vont chercher au bord de la contrainte, là où est le front.
\item \textbf{Le réglage V2 n'est pas optimal partout.} Le poids 25 est le meilleur dans S1, mais dans S2 un poids dix fois plus faible donne 110 simulations au lieu de 142,5.
\end{itemize}

\begin{resultat}
Le succès de la boucle QUBO tient à deux choses qui ne doivent rien au calcul quantique : un modèle affine de la contrainte, optimiste près du front, et un poids de pénalité dans la bonne zone. Ce poids a été choisi en V1 pour une autre contrainte ($C_L\ge0{,}3$) ; qu'il convienne à la V2 n'était pas garanti.
\end{resultat}

\section{Second diagnostic : le krigeage prédit-il mieux ?}
On tire $n$ profils au hasard parmi ceux dont XFOIL a convergé, on entraîne les deux modèles, et on mesure l'erreur sur tous les autres profils (scénario S1, 10 tirages par taille). L'erreur est rapportée à l'écart-type de la grandeur : 1 correspond à un modèle qui prédirait la moyenne partout, 0 à un modèle parfait.

\begin{figure}[htbp]
\centering
% figure générée par tools/make_figures.py : diag_qualite
\begin{tikzpicture}
\begin{groupplot}[group style={group size=3 by 1,horizontal sep=1.0cm},width=0.37\linewidth,height=5.2cm,
  xlabel={profils d'entraînement},xmode=log,log basis x=2,xtick={20,40,80,160},xticklabels={20,40,80,160},
  ymin=0,title style={font=\small},tick label style={font=\scriptsize},label style={font=\small},
  legend style={draw=none,font=\scriptsize,at={(1.75,-0.32)},anchor=north},legend columns=2]
\nextgroupplot[title={traînée $C_D$},ylabel={erreur relative de prédiction}]
\addplot[cblue,solid,very thick,mark=*,mark size=1.6pt] coordinates {(20,0.82251) (40,0.79686) (80,0.78924) (160,0.77945)};\addlegendentry{modèle quadratique}
\addplot[corange,densely dashed,very thick,mark=*,mark size=1.6pt] coordinates {(20,0.8695) (40,0.7926) (80,0.75598) (160,0.76312)};\addlegendentry{krigeage}
\nextgroupplot[title={SER}]
\addplot[cblue,solid,very thick,mark=*,mark size=1.6pt] coordinates {(20,0.13599) (40,0.084809) (80,0.073659) (160,0.070645)};
\addplot[corange,densely dashed,very thick,mark=*,mark size=1.6pt] coordinates {(20,0.11741) (40,0.059505) (80,0.033458) (160,0.018523)};
\nextgroupplot[title={$C_{L,\max}$}]
\addplot[cblue,solid,very thick,mark=*,mark size=1.6pt] coordinates {(20,0.72593) (40,0.6034) (80,0.57003) (160,0.53778)};
\addplot[corange,densely dashed,very thick,mark=*,mark size=1.6pt] coordinates {(20,0.76448) (40,0.64842) (80,0.628) (160,0.68478)};
\end{groupplot}
\end{tikzpicture}

\caption{Erreur de prédiction relative (médiane sur 10 tirages) selon le nombre de profils d'entraînement.}\label{fig:diag-qualite}
\end{figure}

\begin{center}
\begin{tabular}{lcccc}
\toprule
 & \multicolumn{2}{c}{erreur relative, $n=160$} & \multicolumn{2}{c}{erreurs dans $\pm2s$, $n=160$} \\
 & quadratique & krigeage & quadratique & krigeage \\
\midrule
traînée $C_D$ & 0,78 & 0,76 & 98~\% & 91~\% \\
SER & 0,07 & \textbf{0,02} & 95~\% & 94~\% \\
$C_{L,\max}$ & \textbf{0,54} & 0,68 & 96~\% & 85~\% \\
\bottomrule
\end{tabular}
\end{center}

\paragraph{Lecture.}
\begin{itemize}
\item \textbf{La SER} est une fonction lisse : le krigeage la prédit quatre fois mieux. Mais le modèle quadratique la prédit déjà très bien (erreur de 7~\%) : il n'y a presque rien à gagner.
\item \textbf{La traînée} est mal prédite par les deux modèles (erreur de 76 à 78~\%), et ajouter des données n'y change presque rien. L'erreur est dominée par quelques profils à traînée très élevée (décollement) et par les bascules de transition, qu'aucun des deux modèles ne capture.
\item \textbf{$C_{L,\max}$}, la grandeur dont tout dépend, est \emph{moins bien} prédite par le krigeage que par le modèle quadratique, et le krigeage y est trop sûr de lui : 85~\% des erreurs dans $\pm2s$ au lieu de 95~\%.
\end{itemize}

\begin{resultat}
Le krigeage n'est meilleur que là où ce n'est pas utile. Sur les deux grandeurs irrégulières, sa souplesse ne l'aide pas, et son incertitude est sous-estimée, ce qui fausse l'amélioration espérée et la probabilité de faisabilité.
\end{resultat}

\begin{piege}
Ce constat vaut pour le krigeage \emph{tel que réglé ici} : noyau Matérn 5/2, un seul départ pour l'optimisation des hyperparamètres, variables entières traitées comme des réels. Que l'erreur sur $C_{L,\max}$ remonte entre 80 et 160 profils (0,63 puis 0,68) suggère que ce réglage n'est pas robuste. Un krigeage mieux réglé pourrait faire mieux ; on ne l'a pas vérifié.
\end{piege}

\section{Ce que ces diagnostics changent}
\begin{itemize}
\item Ils expliquent pourquoi H1 n'est pas confirmée : sur ce problème, le facteur limitant n'est pas la souplesse du modèle mais la façon d'approcher une contrainte mal prédite.
\item Ils donnent une lecture plus précise de la V2 : la boucle QUBO réussit parce qu'elle combine des modèles simples et un traitement lâche de la contrainte, bien réglé. L'encodage binaire n'y est pour rien : le contrôle sans qubits fait aussi bien. La seule trace de la formulation QUBO dans ce succès est indirecte : c'est elle qui impose le modèle affine.
\item Ils désignent le point faible de toute la chaîne : $C_{L,\max}$, estimé par la non-convergence d'XFOIL sur une rampe d'incidence. C'est l'objet du lot L3 de la fiche projet.
\end{itemize}

\exercices
\begin{exercice}
Pour le profil $(4~\%;\ 0{,}55;\ 12~\%;\ I=2)$ du front, le modèle quadratique prédit $C_{L,\max}=1{,}12\pm0{,}15$ et le krigeage $1{,}04\pm0{,}03$. Calculer les deux probabilités de faisabilité. Lequel des deux modèles laisse une chance à ce profil avec le critère $\mathrm{EI}\times\mathrm{PoF}$ ? Qu'en déduire sur le danger d'une incertitude sous-estimée ?
\end{exercice}
\begin{exercice}
Pourquoi une erreur relative de 0,78 sur la traînée n'empêche-t-elle pas les méthodes de retrouver le front ? (Indication : le front se joue parmi les profils de faible traînée ; que mesure l'erreur quadratique moyenne ?)
\end{exercice}

\begin{pistes}
\begin{itemize}
\item \textbf{Régler automatiquement le poids de la pénalité} au fil de la boucle, par exemple en visant une part cible de profils proposés infaisables, au lieu d'un poids fixe.
\item Entraîner le modèle de $C_{L,\max}$ sur une grandeur plus régulière (le lot L3 propose le maximum de $C_L(\alpha)$ sur une rampe fine) et refaire ce chapitre.
\item Remplacer le modèle de régression de $C_{L,\max}$ par un \emph{classifieur} faisable/infaisable, qui est la seule information dont la sélection a besoin.
\item Mesurer l'erreur des modèles sur les seuls profils proches du front, plutôt que sur toute la grille.
\end{itemize}
\end{pistes}

\chapter{Bilan, limites et pistes}\label{ch:bilan}
\iterbanner{V3}

\section{Réponse à H1}
\begin{resultat}[ : H1 n'est pas confirmée]
Sur le problème V2, l'optimisation bayésienne retrouve le front à chaque fois, mais elle ne bat pas la boucle QUBO : elle demande nettement plus de simulations dans S1 (125 contre 75), et l'écart en sa faveur dans S2 (127,5 contre 142,5) n'est pas significatif. Son modèle plus souple ne l'aide pas.
\end{resultat}

L'ablation dit pourquoi, élément par élément :
\begin{center}
\renewcommand{\arraystretch}{1.25}
\begin{tabularx}{\linewidth}{>{\raggedright\arraybackslash}p{4.2cm}X}
\toprule
Élément & Effet mesuré \\
\midrule
solveur du QUBO (V2) & aucun gain : le solveur exact reste le meilleur \\
encodage binaire & aucun gain : le contrôle sans qubits fait aussi bien dans S1, mieux dans S2 \\
modèle (quadratique ou krigeage) & le krigeage n'améliore aucune des quatre comparaisons \\
sélection (minimum ou EI) & l'amélioration espérée n'améliore aucune des quatre comparaisons \\
traitement de la contrainte de manœuvre & décide du succès : de 0~\% à 100~\% de réussite \\
\bottomrule
\end{tabularx}
\end{center}

\section{Le message pour le jury}
La V2 disait : « le gain vient du modèle de substitution, pas du solveur quantique ». Le lot L1 précise :
\begin{enumerate}
\item le gain vient d'un modèle de substitution \emph{simple} ; un modèle plus riche ne fait pas mieux ;
\item il dépend de la façon dont on approche la contrainte de manœuvre, parce que les meilleurs compromis traînée/SER sont au bord de cette contrainte ;
\item la formulation QUBO n'apporte pas de gain en nombre de simulations. Elle intervient une fois, indirectement : c'est elle qui impose un modèle affine de la contrainte, et ce modèle grossier se trouve être le bon choix ici.
\end{enumerate}
Ce message est moins flatteur pour le quantique que celui de la note d'intention, mais il est défendable : chaque affirmation renvoie à une expérience contrôlée, avec son contrôle et son test statistique. Ce qui reste à la formulation QUBO est ce que la fiche projet annonçait : la mise sous forme de QUBO d'un problème d'ingénierie sous contraintes, la qualité des candidats de QAOA, et l'évolution des ressources avec la taille.

\section{Limites}
\begin{itemize}
\item \textbf{Un seul problème, petit.} 1\,024 profils, deux scénarios radar, des fronts de 4 et 10 profils. Le critère « tout le front exact » est sensible au hasard à cette taille. Rien ne dit que les conclusions tiennent sur la grille de 16\,384 profils du lot L2.
\item \textbf{Un krigeage peu réglé.} Un seul départ pour les hyperparamètres, des variables entières traitées comme des réels, une incertitude sous-estimée sur $C_{L,\max}$. « Le krigeage n'apporte rien » veut dire « ce krigeage-ci, sur ce problème-ci ».
\item \textbf{Une ablation imparfaite sur la sélection.} Les variantes « minimum » et « amélioration espérée » ne traitent pas la contrainte de la même manière ; leur comparaison mêle deux effets.
\item \textbf{Un écart non expliqué.} Dans S2, le contrôle bat la boucle QUBO, et trois différences peuvent l'expliquer.
\item \textbf{Vingt graines, quatorze tests.} Seuls les écarts les plus nets ($p<0{,}002$) sont établis.
\item \textbf{Pas d'amélioration de l'hypervolume.} Le critère multi-objectif de référence n'a pas été essayé.
\end{itemize}

\section{Ce que ce lot change pour la suite de la V3}
\begin{center}
\renewcommand{\arraystretch}{1.25}
\begin{tabularx}{\linewidth}{lX}
\toprule
Lot & Conséquence \\
\midrule
L2 (grande grille) & reprendre l'ablation à l'identique ; \code{bayes\_loop} accepte déjà d'autres tailles de grille. Vérifier d'abord si le poids de pénalité 25 reste dans la bonne zone. \\
L3 (robustesse physique) & devient prioritaire : $C_{L,\max}$ est la grandeur dont tout dépend, et c'est la moins bien prédite. Un $C_{L,\max}$ plus régulier peut changer tout le classement des méthodes. \\
L4 (QAOA bruité) & inchangé. \\
\bottomrule
\end{tabularx}
\end{center}

\begin{pistes}[ : par ordre de priorité]
\begin{enumerate}
\item Expliquer l'écart dans S2 entre le contrôle et la boucle QUBO (un élément à la fois).
\item Refaire le chapitre~\ref{ch:diagnostic} avec le $C_{L,\max}$ du lot L3.
\item Régler le poids de la pénalité au fil de la boucle, plutôt qu'à la main.
\item Compléter l'ablation : « amélioration espérée + pénalité de la boucle QUBO », pour séparer sélection et contrainte.
\item Mieux régler le krigeage (plusieurs départs, noyaux pour variables entières de SMT) avant de conclure définitivement sur le modèle.
\item Coder l'amélioration espérée de l'hypervolume.
\end{enumerate}
\end{pistes}


\appendix
\chapter{Corrigés des exercices}\label{annexe:corriges}

\section*{Chapitre~\ref{ch:question} \iterbadge{V3}}
\textbf{Ex.~1.1.} Sans modèle : tirage aléatoire, NSGA-II, recuit simulé appliqué aux simulations (il ne garde que l'état courant). Avec modèle : la boucle QUBO (polynômes quadratiques des objectifs, modèle affine de la contrainte) et l'optimisation bayésienne (un krigeage par grandeur). Attention : dans la boucle QUBO, le recuit simulé sert à minimiser le \emph{modèle}, pas à choisir directement des simulations.

\textbf{Ex.~1.2.} Les deux méthodes diffèrent par le modèle \emph{et} par la sélection (et par le traitement de la contrainte). Si l'une gagne, on ne sait pas à quel élément l'attribuer. Les deux cases hors diagonale changent un seul élément à la fois ; la case de contrôle vérifie en plus que la réécriture sans QUBO ne change rien.

\section*{Chapitre~\ref{ch:krigeage} \iterbadge{V3}}
\textbf{Ex.~2.1.} $K=k(x_1,x_1)=\sigma^2$, donc $\mu(x)=y_1\,k(x_1,x)/\sigma^2$ et $s^2(x)=\sigma^2-k(x_1,x)^2/\sigma^2$. En $x=x_1$ : $\mu=y_1$ et $s^2=0$, le modèle interpole. Loin de $x_1$ : $\mu\to0$ (la moyenne a priori) et $s^2\to\sigma^2$ (la variance a priori) : une simulation n'apprend rien sur un profil très différent.

\textbf{Ex.~2.2.} $s^2(x)$ ne fait intervenir que $k$, $K$ et $\vect k_\star$, donc les positions des points et les hyperparamètres. C'est une propriété de la loi normale : la variance conditionnelle ne dépend pas de la valeur observée. Ce n'est plus vrai en pratique : les hyperparamètres ($\sigma^2$, $\theta_i$, bruit) sont réglés sur $\vect y$, donc $s$ dépend des valeurs simulées à travers eux. Des données très irrégulières donnent de grands $\theta_i$, donc une incertitude qui remonte vite entre les points.

\textbf{Ex.~2.3.} $\mathrm{Cov}(f(x),f(x'))=\mathbb E\big[\varphi(x)^{\!\top}\vect c\,\vect c^{\top}\varphi(x')\big]=\varphi(x)^{\!\top}\,\mathbb E[\vect c\,\vect c^{\top}]\,\varphi(x')=\lambda^{-1}\varphi(x)^{\!\top}\varphi(x')$, car $\vect c$ est centré de covariance $\lambda^{-1}I$. Une combinaison linéaire de variables normales étant normale, $f$ est bien un processus gaussien.

\section*{Chapitre~\ref{ch:acquisition} \iterbadge{V3}}
\textbf{Ex.~3.1.} Premier profil : $u=(0{,}52-0{,}50)/0{,}02=1$, $\mathrm{EI}=0{,}02\times\Phi(1)+0{,}02\times\varphi_{\mathcal N}(1)=0{,}02\times0{,}841+0{,}02\times0{,}242\approx0{,}0217$. Second : $u=(0{,}52-0{,}60)/0{,}20=-0{,}4$, $\mathrm{EI}=-0{,}08\times\Phi(-0{,}4)+0{,}20\times\varphi_{\mathcal N}(-0{,}4)=-0{,}08\times0{,}345+0{,}20\times0{,}368\approx0{,}0461$. On simule le second, bien que sa prédiction soit moins bonne : son incertitude lui donne une chance de gain plus grande. Le minimum de la prédiction aurait choisi le premier.

\textbf{Ex.~3.2.} En écrivant $\mathrm{EI}=s\,[u\,\Phi(u)+\varphi_{\mathcal N}(u)]$ et en utilisant $\varphi_{\mathcal N}'(u)=-u\,\varphi_{\mathcal N}(u)$, la dérivée du crochet par rapport à $u$ vaut $\Phi(u)$. Comme $\partial u/\partial s=-u/s$ et $\partial u/\partial\mu=-1/s$ : $\partial\mathrm{EI}/\partial s=u\Phi(u)+\varphi_{\mathcal N}(u)-u\Phi(u)=\varphi_{\mathcal N}(u)>0$ et $\partial\mathrm{EI}/\partial\mu=-\Phi(u)<0$. Plus d'incertitude ou une meilleure prédiction augmentent toujours le critère.

\section*{Chapitre~\ref{ch:contraintes} \iterbadge{V3}}
\textbf{Ex.~4.1.} $\mathrm{PoF}=\Phi\big((1{,}15-1{,}2)/0{,}10\big)=\Phi(-0{,}5)\approx0{,}309$. Critère : $0{,}04\times0{,}309\approx0{,}012$. Avec le traitement en dur, $\mu_c=1{,}15<1{,}2$ : le profil est écarté, alors qu'il a près d'une chance sur trois d'être faisable.

\textbf{Ex.~4.2.} Moyenne $0{,}5\times0{,}3+0{,}5\times0{,}5=0{,}4$ ; variance $0{,}25\times0{,}01+0{,}25\times0{,}04=0{,}0125$, écart-type $\approx0{,}112$. L'indépendance est discutable parce que les deux grandeurs dépendent des mêmes paramètres de forme : un nez aigu réduit la SER et modifie la traînée. Si le modèle se trompe sur l'effet du rayon de bord d'attaque, il se trompe sur les deux à la fois. Des erreurs corrélées positivement donneraient une variance plus grande que celle calculée.

\textbf{Ex.~4.3.} Pour $1{,}10$ : $75\times0{,}10^2=0{,}75$, soit les trois quarts de l'étendue d'un objectif normalisé : le profil est pratiquement exclu. Pour $1{,}15$ : $75\times0{,}05^2\approx0{,}19$ : pénalisé, mais il peut encore être choisi si son objectif prédit est bon. La pénalité est donc tolérante sur une marge d'environ 0,05 en dessous de l'exigence.

\section*{Chapitre~\ref{ch:code} \iterbadge{V3}}
\textbf{Ex.~5.1.} Le nombre de simulations dépend beaucoup du tirage initial : certaines graines tombent près du front, d'autres non. Avec les mêmes points initiaux, cette source de variation est commune aux deux méthodes et s'élimine dans la différence graine par graine. On peut alors utiliser un test apparié, plus puissant : il détecte un écart donné avec moins de graines.

\textbf{Ex.~5.2.} $1-\mathrm{PoF}$ est la probabilité que le profil soit infaisable. L'ajouter à un objectif normalisé revient à dire qu'un profil sûrement infaisable « coûte » autant que le pire profil faisable observé, et qu'un profil douteux coûte en proportion de son risque. C'est une pénalité sans paramètre, mais elle se révèle trop sévère sur notre problème (chapitre~\ref{ch:diagnostic}).

\section*{Chapitre~\ref{ch:resultats} \iterbadge{V3}}
\textbf{Ex.~6.1.} Non. L'arrêt réaliste se déclenche quand l'hypervolume stagne, bien avant que le front soit complet : les deux méthodes ont alors la moitié du front. Le critère « front complet » mesure le temps qu'il faut pour trouver les \emph{derniers} profils, et c'est là que les méthodes diffèrent. Pour un usage industriel, où l'on s'arrête sans connaître le front, les deux méthodes sont équivalentes dans S1.

\textbf{Ex.~6.2.} Sous l'hypothèse d'égalité, le nombre de graines où A gagne suit une loi binomiale $\mathcal B(20;\ 0{,}5)$. $\mathbb P(X\ge15)=\big[\binom{20}{15}+\binom{20}{16}+\binom{20}{17}+\binom{20}{18}+\binom{20}{19}+\binom{20}{20}\big]/2^{20}=21\,700/1\,048\,576\approx0{,}021$. En test bilatéral (A ou B nettement meilleur), on double : $0{,}041<0{,}05$. On rejette l'égalité, de justesse.

\section*{Chapitre~\ref{ch:diagnostic} \iterbadge{V3}}
\textbf{Ex.~7.1.} Modèle quadratique : $\mathrm{PoF}=\Phi\big((1{,}12-1{,}2)/0{,}15\big)=\Phi(-0{,}53)\approx0{,}30$. Krigeage : $\Phi\big((1{,}04-1{,}2)/0{,}03\big)=\Phi(-5{,}3)\approx5\cdot10^{-8}$. Avec le modèle quadratique, le profil garde 30~\% de son amélioration espérée et peut être choisi. Avec le krigeage, son critère est pratiquement nul : il ne sera jamais simulé, alors qu'il est faisable ($C_{L,\max}=1{,}20$) et sur le front. Le krigeage se trompe deux fois plus sur la valeur (0,16 contre 0,08), mais ce qui condamne le profil, c'est surtout qu'il annonce une incertitude cinq fois plus petite que celle du modèle quadratique. Une incertitude sous-estimée est plus dangereuse qu'une prédiction fausse.

\textbf{Ex.~7.2.} L'erreur quadratique moyenne est dominée par les grandes erreurs, donc par les quelques profils à traînée très élevée (jusqu'à 319 counts), que le modèle prédit mal. Ces profils ne sont pas sur le front. Parmi les profils de faible traînée, où le front se joue, il suffit que le modèle \emph{classe} à peu près correctement les profils pour que la sélection soit bonne. L'erreur sur toute la grille n'est donc pas la bonne mesure ; d'où la dernière piste du chapitre.


\end{document}
